\chapter{Relativization} \label{Ch3:CH}
\thispagestyle{empty}

The last two chapters were about sets, classes and the relations on them. This one is about formulae: what it takes for one to be true \textit{in} a class rather than in a set, and how much of the universe a formula has to look at before we can tell.

We built $\WF$ in \Cref{Ch1:CH} and saw that Foundation says exactly that every set lies in it. To ask whether the rest of the axioms are true in $\WF$, we need a notion of truth in a class, and satisfaction will not do: it is a notion about sets. Relativization is what replaces it.

With relativization in hand, we ask which formulae mean the same thing in a transitive class as in $V$, and which axioms survive the passage to one. The answers give models of large fragments of $\ZFC$: the levels $V_{\alpha}$, the sets $H_{\kappa}$ of hereditarily small sets, and, by reflection, a transitive model of any finitely many axioms of $\ZFC$.

\section{Relativization to a Class}

Relativization is essentially like satisfaction, just for classes. Our ultimate goal is the following: given $(A, R)$ a class structure and a formula $\psi$, we'd want to define $\psi^{(A, R)}$ to be the relativization (think ``interpretation'') of $\psi$ to $(A, R)$.

Here are some complications with this technique:
\begin{itemize}
    \item $R$ might not be $\in$
    \item There might be parameters involved in the definitions of $A$ and $R$. For instance, if $A$ is a set, then ``$A = \setst{x}{x \in A}$'' takes in a set $A$ as a parameter.
\end{itemize}

Look at $A = \setst{x}{\varphi(x)}$. As a scheme, define $\psi^A\of{\bar{v}}$ (with $\bar{v}$ denoting free variables) by recursion on the complexity of $\psi\of{\bar{v}}$ as follows. We will have atomic, propositional and quantifier steps.

Before starting, rename the variables of $\varphi$ other than $x$ so that none of them occurs in $\psi$. Write $\varphi'$ for the result, so that $A = \setst{x}{\varphi'(x)}$ still, and use $\varphi'$ in place of $\varphi$ throughout. Both directions of the clash matter. The quantifier step below substitutes a variable $u$ of $\psi$ for the free variable $x$, so a $u$ bound in $\varphi$ would be captured and $\varphi\of{u}$ would say something other than ``$u \in A$''; and what that substitution produces is then dropped underneath the quantifiers of $\psi$, so a parameter of $\varphi$ sharing its name with a variable that $\psi$ binds would be captured the other way.
\begin{itemize}
    \item \underline{Atomic:} We interpret equality and membership symbols as equality and membership. Ie,
    \begin{align*}
        \parenth{v_0 = v_1}^A &\text{ is } v_0 = v_1 \\
        \parenth{v_0 \in v_1}^A &\text{ is } v_0 \in v_1 
    \end{align*}

    \item \underline{Propositional:} If $\psi$ is a propositional combination of $\psi_0$ and $\psi_1$, then $\psi^A$ is the same propositional combination of $\psi_0^A$ and $\psi_1^A$.

    \item \underline{Quantifiers:} We relativize a quantifier by restricting it to $A$. Ie,
    \begin{align*}
        \parenth{\forall u \, \psi}^A &\text{ is } \forall u \brac{\varphi'(u) \to \psi^A} \\
        \parenth{\exists u \, \psi}^A &\text{ is } \exists u \brac{\varphi'(u) \land \psi^A}
    \end{align*}
\end{itemize}

The atomic step above is written for the case $R = \in$. For a general $R$ it is that step which changes, because ``$\in$'' is a symbol of the language we are interpreting, and what it is interpreted as in $(A, R)$ is $R$. So if $R = \setst{(x, y)}{\varphi_R\of{x, y, q}}$, then $\parenth{v_0 \in v_1}^{(A, R)}$ is $\varphi_R\of{v_0, v_1, q}$, and since the recursion bottoms out at the atomic subformulae, this replaces every occurrence of $\in$ in $\psi$. The parameter $q$ simply comes along for the ride, which is exactly why $\varphi_R$ needs the same renaming as $\varphi$ did, in both directions: in particular $q$ must not be a variable that $\psi$ binds.

\section{The Relative Consistency of $\ZF$}

Why are we doing all of this? Here's one reason.

\begin{boxtheorem}[working in $\ZFwoF$] \label{Ch3:Thm:ZF_Relativizes_To_WF}
    For each axiom $\sigma \in \ZF$ (\textbf{including} foundation), $\sigma^{\WF}$.
\end{boxtheorem}
The above reads, ``The relativization of each $\ZF$ axiom to $\WF$ is true, working in $\ZFwoF$.'' Another equivalent formulation is that for all axioms $\sigma$ of $\ZF$,
\begin{align}
    \ZFwoF \vdash \sigma^{\WF}
\end{align}

Before drawing the consequence we are after, we isolate the fact about relativization that the argument will run on. We built $\psi^A$ to be the interpretation of $\psi$ in $A$, and an interpretation ought to carry proofs and not just statements: reading every quantifier of a proof as ranging over $A$ turns each of its steps into a step about $A$. The one thing that can go wrong is $A$ being empty, since first-order logic assumes a non-empty domain and proves $\forall u \, \psi \to \exists u \, \psi$; so non-emptiness is carried as a hypothesis rather than as a side condition, being a sentence of the object language like any other.

% [CORRECTED] was: ``if $A$ is a non-empty class (say, with no parameters), then whenever $\sigma \vdash \tau$, $\sigma^A \vdash \tau^A$'' --- read literally that is false, because ``$A$ is non-empty'' is a sentence of the object language rather than a side condition on $A$. Take $\sigma$ to be $\forall u \parenth{u = u}$ and $\tau$ to be $\exists u \parenth{u = u}$, both valid; then $\sigma^A$ is valid too, while $\tau^A$ says exactly that $A \neq \emptyset$, which pure logic does not prove. So the hypothesis belongs on the left of the turnstile. The restriction to sentences is needed for the same reason: with $v$ free, $v = v \vdash \exists u \parenth{u = v}$, but the relativization of the conclusion says $v \in A$, which nothing forces
\begin{boxlemma}[Scheme] \label{Ch3:Lemma:Relativization_Preserves_Proofs}
    Let $A = \setst{x}{\varphi(x)}$, say with no parameters, and let $\sigma$ and $\tau$ be sentences. Then, in the meta-theory, whenever $\sigma \vdash \tau$, we have $\sigma^A, \, \exists x \, \varphi(x) \vdash \tau^A$.
\end{boxlemma}

\Cref{Ch3:Thm:ZF_Relativizes_To_WF} and \Cref{Ch3:Lemma:Relativization_Preserves_Proofs} together yield the following corollary.

\begin{boxcorollary} \label{Ch3:Cor:Con_ZF}
    If $\ZFwoF$ is consistent, then so is $\ZF$.
\end{boxcorollary}
\begin{proof}
    We argue by contraposition. Suppose $\ZF$ is inconsistent. Say $\sigma_1, \ldots, \sigma_n$ are axioms of $\ZFwoF$ and
    \begin{align*}
        \set{\sigma_1, \ldots, \sigma_n} \cup \set{F} \vdash \varphi \land \neg\varphi
    \end{align*}
    Apply \Cref{Ch3:Lemma:Relativization_Preserves_Proofs} with $\sigma$ the conjunction $\sigma_1 \land \cdots \land \sigma_n \land F$ and $\tau$ the contradiction. Relativization commutes with the propositional connectives, by the propositional step of the definition, so what comes out is
    % [CORRECTED] $\exists x \parenth{x \in \WF}$ is added to the left of the display, which had only the relativized axioms --- the lemma's non-emptiness hypothesis is not logically valid and the relativized axioms do not prove it, so it has to be carried here and discharged in the paragraph below, where we are working in $\ZFwoF$ and $\emptyset \in \WF$ is available
    \begin{align*}
        \set{\sigma_1^{\WF}, \ldots, \sigma_n^{\WF}} \cup \set{F^{\WF}} \cup \set{\exists x \parenth{x \in \WF}} \vdash \varphi^{\WF} \land \neg \varphi^{\WF}
    \end{align*}
\end{proof}

% [CORRECTED] ``still under the supposition that $\ZF$ is'' is inserted --- the sentence sits outside the proof environment, whose QED has closed the scope of ``Suppose $\ZF$ is inconsistent'', so as written it asserts the inconsistency of $\ZFwoF$ flatly
Together with the above theorem scheme, we get that $\ZFwoF$ is inconsistent---still under the supposition that $\ZF$ is. Indeed, \Cref{Ch3:Thm:ZF_Relativizes_To_WF} gives $\ZFwoF \vdash \sigma_i^{\WF}$ for each $i$, and also $\ZFwoF \vdash F^{\WF}$; and $\ZFwoF \vdash \exists x \parenth{x \in \WF}$, since $\emptyset \in V_1$. So every hypothesis on the left of that display is a theorem of $\ZFwoF$, and hence $\ZFwoF \vdash \varphi^{\WF} \land \neg\varphi^{\WF}$.

\section{Complexity of Formulae}

Whether a formula says the same thing inside a class as it does in $V$ depends on how its quantifiers range, so we need a way of measuring that. The first step is to give the quantifiers that range over a set, rather than over the whole universe, notation of their own.

\begin{boxnotation}[Bounded Quantifiers] \hfill
    \begin{itemize}
        \item $\forall u \in v \, \varphi$ stands for $\forall u \brac{u \in v \to \varphi}$

        \item $\exists u \in v \, \varphi$ stands for $\exists u \brac{u \in v \land \varphi}$
    \end{itemize}
    Quantifiers written this way are said to be \textbf{bounded}.
\end{boxnotation}

A formula all of whose quantifiers are bounded can only ever look inside the sets it is handed, and such formulae sit at the bottom of the hierarchy we are about to build.

\begin{boxdefinition}[$\Delta_0$ Formulae]
    A formula is $\Delta_0$ iff every quantifier occurring in it is bounded.
\end{boxdefinition}

Everything above the bottom level is obtained by putting a block of unbounded quantifiers, all of the same kind, in front of a $\Delta_0$ formula.

\begin{boxdefinition}[$\Sigma_1$ and $\Pi_1$ Formulae] \hfill
    \begin{itemize}
        \item $\Sigma_1$ formulae look like
        \begin{align*}
            \exists u_0 \exists u_1 \cdots \exists u_n \, \varphi
        \end{align*}
        where $\varphi$ is $\Delta_0$.

        \item $\Pi_1$ formulae look like
        \begin{align*}
            \forall u_0 \forall u_1 \cdots \forall u_n \, \psi
        \end{align*}
        where $\psi$ is $\Delta_0$.
    \end{itemize}
\end{boxdefinition}

There are no $\Delta_1$ formulae. On the shapes above, a formula cannot begin with both an unbounded $\exists$ and an unbounded $\forall$; and allowing either block to be empty only puts the $\Delta_0$ formulae into both classes, which buys us nothing we did not already have. What there are is $\Delta_1^T$ formulae, and to get them we first have to fix a theory to do the translating for us.

\begin{boxdefinition}[Complexity Relative to a Theory]
    Let $T$ be a theory and let $\psi\of{\bar{u}}$ be a formula. We say $\psi\of{\bar{u}}$ is $\Delta_0^T$ iff there is a $\Delta_0$ formula $\varphi\of{\bar{u}}$ such that
    \begin{align*}
        T \vdash \forall \bar{u} \brac{\psi\of{\bar{u}} \leftrightarrow \varphi\of{\bar{u}}}
    \end{align*}
    Similarly for $\Sigma_1^T$ and $\Pi_1^T$. We say $\psi\of{\bar{u}}$ is $\Delta_1^T$ iff it is both $\Sigma_1^T$ and $\Pi_1^T$.
\end{boxdefinition}

Keep going to get $\Sigma_n$ and $\Pi_n$ formulae, and also $\Sigma_n^T$, $\Pi_n^T$ and $\Delta_n^T$ formulae.

\section{Absoluteness} \label{Ch3:Sec:Absoluteness}

\subsection{Absoluteness Theorems}

We will constantly be comparing what a formula says in a smaller class with what it says in a bigger one, so we give the possible outcomes names.

\begin{boxdefinition}[Absoluteness]
    Let $M \ssq N$ be classes and let $\vp\of{\bar{u}}$ be a formula. We say $\vp\of{\bar{u}}$ is
    \begin{itemize}
        \item \textbf{downwards absolute} from $N$ to $M$ iff for all $\bar{x} \in M^{<\omega}$ matching $\bar{u}$, $\vp^N\of{\bar{x}} \to \vp^M\of{\bar{x}}$.
        \item \textbf{upwards absolute} from $M$ to $N$ iff for all $\bar{x} \in M^{<\omega}$ matching $\bar{u}$, $\vp^M\of{\bar{x}} \to \vp^N\of{\bar{x}}$.
        \item \textbf{absolute} between $M$ and $N$ iff it is both, ie, iff for all $\bar{x} \in M^{<\omega}$ matching $\bar{u}$, $\vp^M\of{\bar{x}} \lr \vp^N\of{\bar{x}}$.
    \end{itemize}
    When $N = V$, we just say $\vp\of{\bar{u}}$ is downwards absolute, upwards absolute or absolute \textbf{for $M$}.
\end{boxdefinition}

We will want to relativize classes as well as formulae, and it is convenient to have notation for it.

\begin{boxnotation}
    Say $M$ is a transitive class and $A$ another class. Say $A = \setst{x}{\vp(x)}$. We might write $A^M$ to mean $\setst{x \in M}{\vp^M(x)}$ even though it may depend on the choice of $\varphi$ with this notation.

    $A^M$ is what $M$ thinks $A$ is. When it agrees with $A \cap M$, ie, when $\vp(x)$ is absolute for $M$, we say $A$ is \textbf{absolute} for $M$.
\end{boxnotation}

We will use this constantly: in \Cref{Ch3:Subsec:Recursion} we ask whether a class function $G$ defined by recursion satisfies $G^M = G \restriction \parenth{A \cap M}$, and in \Cref{Ch3:Subsec:Ordinals} whether $\OR^M = \OR \cap M$.

Every absoluteness theorem we prove will have an assumption that the smaller model (or sometimes, both models) are transitive. There are no absoluteness theorems in this course that do not mention transitivity, and it is a common mistake to try and apply them in cases where there isn't transitivity.

\begin{boxlemma}[$\Delta_0$-Absoluteness]
    Let $M$ and $N$ be classes\footnote{possibly proper classes, possibly sets}, with $M \ssq N$. Assume $M$ is transitive. Then, for every $\Delta_0$ formula $\varphi\of{\bar{u}}$ and $\bar{x} \in M^{<\omega}$ (matching $\bar{u}$ in length),
    \begin{align*}
        \varphi^M\of{\bar{x}} \lr \varphi^N\of{\bar{x}}
    \end{align*}
    That is, $\varphi\of{\bar{u}}$ is absolute between $M$ and $N$.
\end{boxlemma}
\begin{proof}
    We argue by induction on the complexity of $\varphi\of{\bar{u}}$.

    \begin{itemize}
        \item \underline{Atomic Case.} Note that
        \begin{align*}
            \parenth{u_0 = u_1}^M
            \qquad
            u_0 = u_1
            \qquad
            \parenth{u_0 = u_1}^N
        \end{align*}
        are all exactly the same. Similarly,
        \begin{align*}
            \parenth{u_0 \in u_1}^M
            \qquad
            u_0 \in u_1
            \qquad
            \parenth{u_0 \in u_1}^N
        \end{align*}
        are all exactly the same.

        \item \underline{Propositional Case.} If $\varphi$ is some propositional combination of $\varphi_0$ and $\varphi_1$, then the statement for $\vp$ is clear from the induction hypotheses for $\vp_0$ and $\vp_1$.

        \item \underline{Bounded Quantifiers.} Say $\vp\of{\bar{u}}$ is of the form $\exists v \in w \, \parenth{\psi\of{\bar{u}, v, w}}$. (This is enough, because bounded universal can be formulated as a negation of a bounded existential, and the propositional case covers negations.)

        Let $\bar{x} \in M^{<\omega}$ and $z \in M$. Then, the following are equivalent:
        \begin{itemize}
            \item $\brac{\exists v \in z, \, \psi\of{\bar{x}, v, z}}^M$
            \item There is $y \in z \cap M$ such that $\psi^M\of{\bar{x}, y, z}$
            \item There is $y \in z$ such that $\psi^M\of{\bar{x}, y, z}$ (because $M$ is transitive, so $z = z \cap M$)
            \item There is $y \in z \cap N$ such that $\psi^M\of{\bar{x}, y, z}$ (because $z \ssq M \ssq N$, so $z = z \cap N$)
            \item There is $y \in z \cap N$ such that $\psi^N\of{\bar{x}, y, z}$
            \item $\brac{\exists v \in z, \, \psi\of{\bar{x}, v, z}}^N$
        \end{itemize}
    \end{itemize}
\end{proof}

One level up, we only get one direction each.

\begin{boxlemma}
    In the same situation as above, we have upward absoluteness from $M$ to $N$ for $\Sigma_1$ formulae and downward absoluteness from $N$ to $M$ for $\Pi_1$ formulae.
\end{boxlemma}

Relative to a theory, the same holds for formulae that the theory proves equivalent to these shapes.

\begin{boxlemma}
    Let $T$ be a theory and $M \ssq N$ be models of $T$ (with $M$ transitive). Then,
    \begin{itemize}
        \item $\Sigma_1^T$ formulae are upwards-absolute
        \item $\Pi_1^T$ formulae are downwards-absolute
        \item $\Delta_1^T$ formulae are absolute
    \end{itemize}
\end{boxlemma}

\subsection{Finite Sets}

For the rest of this section, let $M$ be a transitive class model of $\ZFwoP$. We start with the finite sets.

\begin{boxlemma}
    $[M]^{< \omega} \ssq M$.
\end{boxlemma}
\begin{proof}
    Given $x_1, \ldots, x_n \in M$, $\set{x_1, \ldots, x_n} = \bigcup_{j=1}^{n} \set{x_j}$.
\end{proof}

The formula ``$A$ is finite'' is absolute for $M$. In fact, we can do better.

\begin{boxlemma}
    For each $n < \omega$, the formula ``$\abs{A} = n$'' is absolute for $M$.
\end{boxlemma}
\begin{proof}
    \hfill 
    \begin{description}
        \item[\underline{Upwards absoluteness.}] Say $f \in M$ and
        \begin{align*}
            \parenth{f : n \iso A}^M
        \end{align*}
        where $f : n \iso A$ means $f$ is injective and surjective.

        We know that $f : n \iso A$ because we already know it is absolute.

        \item[\underline{Downwards absoluteness.}] We have that $f : n \iso A$ in $V$. But $f \ssq_{\text{finite}} n \times A \ssq M$. And indeed we've checked already that Cartesian products are absolute.

        Then, by the previous lemma, $f \in M$.

        Now again, by $\Delta_0$-absoluteness,
        \begin{align*}
            \parenth{f : n \iso A}^M
        \end{align*}
    \end{description}
\end{proof}

Being finite and having a given finite size are the same thing, so we get the following for free.

\begin{boxlemma}
    ``$A$ infinite'' is absolute for $M$.
\end{boxlemma}
\begin{proof}
    Clear from previous.
\end{proof}

The power set of a finite set is also computed correctly.

\begin{boxlemma}
    If $A \in M$ and $A$ is finite (in both $M$ and $V$), then $\Powset(A) \in M$. Therefore, $A \mapsto \Powset(A)$, restricted to finite sets $A$, is an absolute class function for $M$. % [CORRECTED] ``restricted to finite sets $A$'' inserted --- for infinite $A$ it fails, eg, $\Powset(\omega) \notin \HC$

    For the remainder of this argument, let's work in $M$.
    
    We argue by recursion. We have that $A$ is finite, so there is some $n < \omega$ and $f : n \iso A$. Denote by $2^n$ the set of all $n$-tuples of $0$s and $1$s. For each $x \in 2^n$, write
    \begin{align*}
        B_x = \setst{f(i)}{x_i = 1} \ssq A
    \end{align*}
    In particular, we have a mapping $x \mapsto B_x$. Finally, we can conclude that
    \begin{align*}
        \Powset(A) = \setst{B_x}{x \in 2^n}
    \end{align*}
\end{boxlemma}

\subsection{Absoluteness of Recursive Definitions} \label{Ch3:Subsec:Recursion}

Often, we have $(A, R)$ well-founded and set-like, and some $F : A \times V \to V$. We then define $G : A \to V$ by recursion so that
\begin{align*}
    \forall y \in A, \, G(y) = F\of{y, G \restriction \setst{x}{x \, R \, y}}
\end{align*}
Now, we'd like to know that if the ``data'' above is absolute, then so is $G$. How do we show this?

Ok, let's list our assumptions (recalling that a class function is absolute iff the defining formula is absolute, and also $M$ is closed under the values):
\begin{itemize}
    \item $A$ and $R$ are absolute
    \item $F$ is an absolute class function
    \item $y \mapsto \setst{x}{x \, R \, y}$ is an absolute class function
\end{itemize}
We can conclude from this the following two lemmas.

\begin{boxlemma}
    $\brac{\text{$(A, R)$ is set-like and well-founded}}^M$
\end{boxlemma}
\begin{proof}[Proof sketch]
    Note that $(A, R)$ is well-founded iff
    \begin{align*}
        \forall S, \, \brac{\emptyset \ne S \ssq A \to \exists x \in S, \, \forall y \in S,\, \neg \parenth{y \, R \, x}}
    \end{align*}
    We can essentially rewrite the expression inside brackets using quantifiers that are bounded over $(A, R)$ (subject to the universal $\forall S$). The bracketed expression is absolute between $M$ and $V$, and the universal out front makes it downwards absolute from $V$ to $M$.

    Note that set-like is downwards absolute from $V$ to $M$ because we assumed that $y \mapsto \setst{x}{x \, R \, y}$ and $(A, R)$ are all absolute.
\end{proof}

With that, $G$ is computed correctly in $M$.

\begin{boxlemma}
    $G$ is absolute for $M$. In particular, for every $y \in A^M$,
    \begin{align*}
        G^M(y) = F^M\of{y, G^M \restriction \setst{x}{x \, R^M y}}
    \end{align*}
    and
    \begin{align*}
        G^M = G \restriction A^M = G \restriction \parenth{A \cap M}
    \end{align*}
\end{boxlemma}
\begin{proof}[Proof sketch]
    Just show that for all $y$ in $A^M = A \cap M$, we have $G^M(y) = G(y)$ by induction on $y$.
\end{proof}

There are many useful applications of absoluteness.

\begin{boxexample}[Defining Ordinal Arithmetic]
    There are many ways we can use absoluteness to define ordinal arithmetic.
    \begin{itemize}
        \item \underline{Defining it recursively.} For example, we can define $\alpha + 0 = \alpha$, $\alpha + \parenth{\beta + 1} = \parenth{\alpha + \beta} + 1$, and if $\gamma$ is a limit ordinal, we can define $\alpha + \gamma := \sup_{\beta < \gamma}\of{\alpha + \beta}$.

        \item \underline{Defining it using Mostowski Collapse of well-orderings.} For example, we can define
        \begin{align*}
            \alpha + \beta :=
            \type\of{\parenth{\parenth{\set{0} \times \alpha} \cup \parenth{\set{1} \times \beta}, <_{\lex}}}
        \end{align*}
        Note that $\type$ is the range of the Mostowski Collapse. Indeed,
        \begin{align*}
            G(y) = \setst{G(x)}{x <_{\lex} y}
        \end{align*}
        is absolute.
    \end{itemize}
\end{boxexample}

\subsection{Ordinals in Transitive Models} \label{Ch3:Subsec:Ordinals}

% One final remark for today. We know that $M$ is a transitive class. If $M$ is a proper class, then $\OR^M = \OR$. If $M$ is a set, then $\OR^M \in \OR$ (ie, if $M$ is a set, then as a set of ordinals that is transitive, it must itself be a 

If $M$ is a transitive model of $\ZFwoP$, we claim that the statement ``$\alpha$ is an ordinal'' is absolute for $M$. % [CORRECTED] was: ``$\alpha$ is absolute for $M$'' --- absoluteness is a property of formulae, and what the next display unpacks is the formula ``$\alpha$ is an ordinal'\,'

But what does it even mean? It means
\begin{align*}
    \OR^M = \OR \cap M
\end{align*}
As $M$ is transitive, $\OR^M$ is transitive, so either $\OR^M \in \OR$ or $\OR^M = \OR$. It turns out:
\begin{enumerate}
    \item $\OR^M \in \OR$ iff $M$ is a set
    \item $\OR^M = \OR$ iff $M$ is a proper class
\end{enumerate}
To see this, note that foundation is equivalent to $(V, \in)$ being well-founded. It's clear that $(V, \in)$ is set-like: if $y$ is a set, then everything that is $\in$-below $y$ is just $y$, which is a set. So by the Theorem Scheme on Recursion (which is true in $\ZFwoP$), we have a class function $G : V \to \OR$ such that $\forall y$, $G(y) = \sup_{x \in y}\of{G(x) + 1}$.

Now if we also give ourselves the power set axiom, then we can define the $V$-hierarchy, and conclude that $G(y) = \rank{y}$. But without power sets, \textit{we cannot even define the $V$-hierarchy}, so there is no meaningful notion of rank!

By the theorem scheme on absoluteness of recursively defined class functions, $G$ is an absolute class function for $M$. In other words, $G^M = G \restriction M$. Also, $G^M : M \to \OR^M$ is surjective. % [FILLED] the surjectivity and the two cases of the dichotomy are argued here; the lecture asserted them
Indeed, each value of $G^M$ is an ordinal in $M$, and $G(\alpha) = \alpha$ for every ordinal $\alpha$, by induction: $G(\alpha) = \sup_{\beta < \alpha}\of{G(\beta) + 1} = \sup_{\beta < \alpha}\of{\beta + 1} = \alpha$. So every $\alpha \in \OR^M = \OR \cap M$ is $G^M(\alpha)$.

This gives both halves of the dichotomy.
\begin{itemize}
    \item If $M$ is a set, then $\OR \cap M = G[M]$ is a set by Replacement, and a transitive set of ordinals is an ordinal, so $\OR^M \in \OR$.
    \item If $M$ is a proper class, suppose $\OR \cap M$ were a set, hence an ordinal $\alpha$. In $V$ we do have Power Set, so $G(y) = \rank{y}$ for every $y$. So, for every $y \in M$, $\rank{y} = G^M(y) \in \OR \cap M = \alpha$, ie, $y \in V_{\alpha}$. Then $M \ssq V_{\alpha}$, and $M$ would be a set by Comprehension. So $\OR^M$ is a transitive proper class of ordinals, which can only be $\OR$.
\end{itemize}

Now let's study the uncountable ordinals in $M$.

Denote by $\omega_1^M$ what ``$M$ thinks'' is the least uncountable ordinal. If $\parenth{\omega_1 \text{ exists}}^M$, ie, if the notion of the existence of the least uncountable ordinal is even definable in $M$, then $\omega_1^M \leq \omega_1$ (where $\le$ is just the usual $\le$ on ordinals---recall that $\OR$ is absolute, which means precisely that whatever a class \textit{thinks} is an ordinal is, indeed, an ordinal in $V$).

\section{Relativizing the Axioms of $\ZFC$}

We now put absoluteness to work on the axioms themselves, asking which of them hold relativized to which transitive classes. This is what \Cref{Ch3:Thm:ZF_Relativizes_To_WF} needed for $\WF$, and it also finds models of large fragments of $\ZFC$ among the levels of the $V$-hierarchy.

\subsection{The First Few Axioms}

For this discussion, fix a class $M$. Let's list out the axioms of $\ZFwoF$, relativized to $M$.

We start with Extensionality, which relativized to $M$ reads
\begin{align*}
    \brac{\forall A \forall B \brac{\forall x \parenth{x \in A \lr x \in B} \to A = B}}^M
\end{align*}
Or, equivalently,
\begin{align*}
    \brac{\forall A \forall B \brac{\parenth{A \ssq B \land B \ssq A} \to A = B}}^M
\end{align*}
where $A \ssq B$ is the $\Delta_0$ formula $\parenth{\forall x \in A}\parenth{x \in B}$. Indeed, all sub-formulae of the above are $\Delta_0$.

This is all we need, as long as $M$ is transitive.

\begin{boxlemma}[Extensionality]
    If $M$ is transitive, then $(\text{Extensionality})^M$.
\end{boxlemma}
\begin{proof}
    By $\Delta_0$-absoluteness, for $A, B \in M$ the formula $\parenth{A \ssq B \land B \ssq A} \to A = B$ says the same thing in $M$ as it does in $V$, where it is true. So it holds in $M$ for all $A, B \in M$, which is exactly $(\text{Extensionality})^M$.
\end{proof}

From now on, assume that \underline{$M$ is transitive}, so that $(\text{Extensionality})^M$.

Comprehension needs $M$ to contain enough subsets of its members.

\begin{boxlemma}[Comprehension]
    If $M$ is closed under subsets, ie, $A \in M$ and $B \ssq A$ imply $B \in M$, then $(\text{Comprehension})^M$.
\end{boxlemma}
\begin{proof} % [FILLED] the analysis of Comprehension^M, and the closure condition that secures it, are supplied here; the lecture broke off after the transitivity assumption
    For each formula $\vp\of{x, \bar{y}}$, the corresponding instance relativized to $M$ is
    \begin{align*}
        \forall A \in M \, \forall \bar{y} \in M \, \exists B \in M \, \forall x \in M \brac{x \in B \lr \parenth{x \in A \land \vp^M\of{x, \bar{y}}}}
    \end{align*}
    As $M$ is transitive, any such $B$ and the given $A$ are subsets of $M$, so the bracket pins $B$ down completely: it must be $\setst{x \in A}{\vp^M\of{x, \bar{y}}}$. So Comprehension relativizes to $M$ iff each of these sets lies in $M$, and in particular it does whenever $M$ is closed under subsets.
\end{proof}

From now on, assume $(\text{Comprehension})^M$ as well.

Pairing only needs $M$ to be closed under pairs.

\begin{boxlemma}[Pairing]
    If $\set{x, y} \in M$ for all $x, y \in M$, then $(\text{Pairing})^M$.
\end{boxlemma}
\begin{proof} % [FILLED] the lecture left Pairing blank
    Relativized to $M$, Pairing says
    \begin{align*}
        \forall x \in M \, \forall y \in M \, \exists A \in M \, \parenth{x \in A \land y \in A}
    \end{align*}
    If $\set{x, y} \in M$, then it is a witness. (Here we did not even need transitivity, because the formula under the quantifiers is atomic.)
\end{proof}

\subsection{The Levels of the $V$-Hierarchy}

Let's see which axioms hold in the levels of the $V$-hierarchy.

\begin{boxproposition}
    For each axiom $\sigma$ of $\set{\text{Extensionality}, \text{Comprehension}, \text{Pairing}, \text{Unions}, \text{Power Sets}}$, for all limit ordinals $\alpha$, $\sigma^{V_{\alpha}}$.
\end{boxproposition}

At $V_{\omega}$, Replacement holds as well.

\begin{boxproposition}
    For each axiom $\sigma$ of $\set{\text{Extensionality}, \text{Comprehension}, \text{Pairing}, \text{Unions}, \text{Power Sets}, \text{Replacement}}$, $\sigma^{V_{\omega}}$.
\end{boxproposition}

Above $V_{\omega}$, a limit level need not satisfy Replacement.

\begin{boxnexample}[$V_{\omega + \omega}$]
    We can show that Replacement does not relativize to $V_{\omega + \omega}$. Indeed, consider the class function $G : n \mapsto \omega + n$. It's definable in two ways:
    \begin{enumerate}
        % [CORRECTED] was: $\exists g \brac{\dom(g) = n + 1 \land g(0) = w \land g(i + 1) = g(i) \cup \set{g(i) = g(i) + 1 \land g(n) = S}}$ --- $i$ was never quantified, the brace closed around $g(n) = S$, $w$ stood for $\omega$, and nothing said $g$ is a function, so it did not express $S = \omega + n$
        \item $G(n) = S$ iff $\exists g \brac{g \text{ is a function} \land \dom{g} = n + 1 \land g(0) = \omega \land \forall i < n \, \parenth{g(i + 1) = g(i) \cup \set{g(i)}} \land g(n) = S}$
        % [FILLED] the second definition and the argument after it are supplied here; the lecture announced two definitions, gave one, and stopped
        \item $G(n) = S$ iff $\forall g \brac{\parenth{g \text{ is a function} \land \dom{g} = n + 1 \land g(0) = \omega \land \forall i < n \, \parenth{g(i + 1) = g(i) \cup \set{g(i)}}} \to g(n) = S}$
    \end{enumerate}
    Both say that $S$ is the last value of the function $g$ on $n + 1$ that starts at $\omega$ and takes a successor at each step, and in both, everything after the unbounded quantifier over $g$ is $\Delta_0$ in $g$, $n$, $S$ and the parameter $\omega$. So the first is $\Sigma_1$ and the second is $\Pi_1$.

    Now fix $n < \omega$. The unique $g$ in question is $\setst{\parenth{i, \omega + i}}{i \leq n}$, and each of its pairs lies in $V_{\omega + n + 3}$, so $g \in V_{\omega + n + 4} \ssq V_{\omega + \omega}$. Hence, for $S \in V_{\omega + \omega}$, the first formula relativized to $V_{\omega + \omega}$ holds iff $S = \omega + n$: the witness $g$ is available inside $V_{\omega + \omega}$, and the $\Delta_0$ part means the same thing there as in $V$ by $\Delta_0$-absoluteness. (The second formula relativizes correctly too, by downward absoluteness of $\Pi_1$ formulae, so it does not matter which definition we use.)

    So the hypothesis of the instance of Replacement for this formula, with $A = \omega \in V_{\omega + \omega}$ and parameter $\omega$, holds in $V_{\omega + \omega}$. If the conclusion did too, we would have $B \in V_{\omega + \omega}$ with $\omega + n \in B$ for every $n < \omega$. But $B \in V_{\omega + k + 1}$ for some $k < \omega$, so $B \ssq V_{\omega + k}$, and then $\omega + k \in V_{\omega + k}$, ie, $\rank{\omega + k} < \omega + k$, which is absurd since ordinals are their own ranks.
\end{boxnexample}

We now continue going through the axioms.

Let's look at Foundation.

We want to see if $\F^M$, that is,
\begin{align*}
    \brac{\forall A \parenth{A \ne \emptyset \to \exists x \in A, \, \forall y \in A, \, y \notin x}}^M
\end{align*}
Consider the statement
\begin{align*}
    \forall A \in M \parenth{A \ne \emptyset \to \exists x \in A, \, \forall y \in A, \, y \notin x}
\end{align*}
This is true if $M \ssq \WF$.

We know the following:
\begin{itemize}
    \item if $A \in \WF$, then $A \subset \WF$
    \item if $x \in A$ with $\rank{x} = \min\setst{\rank{y}}{y \in A}$ then $x$ is $\in$-minimal in $A$.
    \item if $y \in x$ then $\rank{y} < \rank{x}$
\end{itemize}
In particular, $\F^{\WF}$ and $\F^{V_{\alpha}}$.

Finally, for infinity, we can show that if $\omega \in M$, then infinity relativizes to $M$; in particular, infinity relativizes to $\WF$ and to $V_{\alpha}$ for $\alpha > \omega$ a limit.

We're now done showing that every axiom of $\ZF$ without infinity relativizes to $V_{\omega}$, and furthermore, that every axiom of $\ZF$ relativizes to $\WF$. As we saw earlier, this is what gives us \Cref{Ch3:Cor:Con_ZF}: if $\ZFwoF$ is consistent, then so is $\ZF$.

\subsection{Choice}

Now, let's study choice. We know that under $\ZFwoF$, this is equivalent to the well-ordering principle.

We want to show, effectively, that
\begin{align*}
    \brac{\forall A \exists R \parenth{(A, R) \text{ is a w.o.}}}^M
\end{align*}
% [CORRECTED] was: ``This is equivalent to $\forall A \in M, \, \exists R \in M \, \brac{\parenth{(A, R) \text{ is a w.o.}}}^M$ because \ldots'' --- that equivalence holds by the definition of relativization, so the reason given did no work; the absoluteness is what transfers ``w.o.'' from $V$ down to $M$, which is how the argument below uses it
This follows from
\begin{align*}
    \forall A \in M, \, \exists R \in M \, \parenth{(A, R) \text{ is a w.o.}}
\end{align*}
because the statement ``$(A, R)$ is a w.o.'' is $\Pi_1$, and thus, downwards absolute.

Indeed, the reason this statement is $\Pi_1$ is that $(A, R)$ is a w.o. iff $(A, R)$ is a l.o. (which is a $\Delta_0$ statement) and $(A, R)$ is well-founded (which is $\Pi_1$ because it says $\forall S \brac{\parenth{\emptyset \ne S \ssq A} \to \exists x \in S \, \forall y \in S \, \neg\parenth{y \, R \, x}}$, and everything after $\forall S$ is $\Delta_0$: the quantifiers over $x$ and $y$ are bounded by $S$, and $\neg\parenth{y \, R \, x}$ says that the ordered pair $\parenth{y, x}$ is not a member of $R$, which needs only quantifiers bounded by $R$). % [FILLED] the reason well-foundedness is Pi_1 is supplied here

In the background, now, assume $\ZFCwoF$. Let $A \in \WF$ and suppose $A \ssq V_{\alpha}$ for some $\alpha$. In $V$, we have some $R \ssq A \times A$ such that $R$ is a w.o. Then $R \ssq V_{\alpha} \times V_{\alpha}$, so $R \in V_{\alpha + 3} \subset \WF$. By $\Pi_1$ downwards absoluteness, we conclude that $\brac{(A, R) \text{ is a w.o.}}^{\WF}$.

\begin{boxcorollary}
    $(\mathbf{AC})^{\WF}$, $(\mathbf{AC})^{\V_{\alpha}}$ for all limit $\alpha$.
\end{boxcorollary}

Since $\WF$ satisfies Choice, the relative consistency result of \Cref{Ch3:Cor:Con_ZF} extends to $\ZFC$.

\begin{boxcorollary}
    If $\ZFCwoF$ is consistent, then so is $\ZFC$.
\end{boxcorollary}

Being a well-ordering is in fact also $\Sigma_1$, over a weak enough theory, so it is absolute for transitive models of that theory and not merely downwards absolute.

\begin{boxlemma}
    % [CORRECTED] was: ``without power sets, infinity and foundation'' --- the proof needs Foundation to make ``$x$ is an ordinal'' a $\Delta_0$ condition: without it, a transitive, $\in$-linear, ill-founded set can witness the displayed $\Sigma_1$ form for a relation that is not a well-ordering. That only shows this particular $\Sigma_1$ form fails; that no $\Sigma_1$ form works without Foundation was argued by the adjudicator via Boffa's anti-foundation axiom, and is not checked here
    If we denote by $T$ the theory consisting of $\ZF$ without power sets and infinity, the statement ``$(A, R)$ is a w.o.'' is also $\Sigma_1^T$.
\end{boxlemma}
\begin{proof}
    $(A, R)$ is a w.o. iff there is some $F : A \to \OR$ such that $\forall x, y \in A \parenth{x \, R \, y \to F(x) < F(y)}$ % [FILLED] the rest of the proof is supplied here; the lecture stated the characterization and stopped
    and $(A, R)$ is a l.o. From left to right, take $F$ to be the collapsing map of \Cref{Ch1:Thm:Mostowski_WO}, which is order-preserving with ordinal values; its construction is a recursion along a well-ordering of a set, and uses Replacement but neither Power Set nor Infinity. From right to left, given a nonempty $S \ssq A$, the set $F[S]$ is a nonempty set of ordinals, so it has a least element $F(x)$ with $x \in S$, and no $y \in S$ has $y \, R \, x$, as that would give $F(y) < F(x)$.

    So, in $T$, ``$(A, R)$ is a w.o.'' is equivalent to
    \begin{align*}
        \text{$(A, R)$ is a l.o.} \land \exists F \brac{F \text{ is a function} \land \dom{F} = A \land \forall x \in A \, \parenth{F(x) \in \OR} \land \forall x, y \in A \, \parenth{x \, R \, y \to F(x) \in F(y)}}
    \end{align*}
    Everything after $\exists F$ is $\Delta_0$ provided ``$F(x) \in \OR$'' is, and this is where Foundation is used: in its presence, an ordinal is just a transitive set that is linearly ordered by $\in$, which is a $\Delta_0$ condition.
\end{proof}

\subsection{Inaccessible Cardinals}

To get a level of the $V$-hierarchy satisfying all of $\ZFC$, we need a level that is closed under everything, and that calls for some definitions about cardinals.

\begin{boxdefinition}[(Strong) Limit Cardinals]
    $\mu$ is a
    \begin{itemize}
        \item \textbf{limit cardinal} iff for all cardinals $\kappa < \mu$, $\kappa^+ < \mu$.
        \item \textbf{strong limit cardinal} iff for all cardinals $\kappa < \mu$, $2^{\kappa} < \mu$.
    \end{itemize}
\end{boxdefinition}

A cardinal can be a limit while still being reachable in fewer steps than its size, which is what cofinality measures.

\begin{boxdefinition}[Cofinality]
    % [CORRECTED] was: ``of a set $A$'' --- ``unbounded in $A$'' needs an order on $A$, and the notes only take cofinalities of ordinals
    The \textbf{cofinality} of an ordinal $A$, denoted $\cf(A)$, is the least $\abs{B}$ such that there is a function $f : B \to A$ whose image $f[B]$ is unbounded in $A$.
\end{boxdefinition}

The cardinals that cannot be reached in fewer steps get a name.

\begin{boxdefinition}[Regularity/Singularity]
    A cardinal $\mu$ is \textbf{regular} if $\cf(\mu) = \mu$. If $\mu$ is not regular, we say $\mu$ is \textbf{singular}.
\end{boxdefinition}

Putting these together gives the cardinals we are after.

\begin{boxdefinition}[Weak/Strong Inaccessibility]
    We say a cardinal $\mu$ is
    \begin{itemize}
        \item \textbf{weakly inaccessible} iff $\mu$ is an uncountable, regular limit cardinal.
        \item \textbf{strongly inaccessible} iff $\mu$ is an uncountable, regular, strong limit cardinal.
    \end{itemize}
\end{boxdefinition}

A strongly inaccessible cardinal is a fixed point of the $\beth$ function.

\begin{boxproposition} \label{Ch3:Prop:Inaccessible_Beth}
    If $\mu$ is strongly inaccessible, then $\beth_{\mu} = \mu$.
\end{boxproposition}
\begin{proof}
    By induction, $\beth_{\alpha} \geq \alpha$ for every ordinal $\alpha$, so $\beth_{\mu} \geq \mu$.

    It's enough to show that $\forall \alpha < \mu$, $\beth_{\alpha} < \mu$, because the fact that $\mu$ is a limit ordinal means that, by definition, $\beth_{\mu} = \sup_{\alpha < \mu} \beth_{\alpha}$.

    Now perform induction on $\alpha < \mu$.
    
    Base case: $\beth_0 < \mu$ because $\beth_0 = \omega$ and $\mu$ is uncountable.

    Successor case: if $\alpha + 1 < \mu$ and $\beth_{\alpha} < \mu$, then $\beth_{\alpha + 1} = 2^{\beth_{\alpha}} < \mu$, because $\beth_{\alpha}$ is a cardinal below $\mu$ and $\mu$ is a strong limit.

    Limit case: if $\alpha < \mu$ is a limit and $\beth_{\gamma} < \mu$ for all $\gamma < \alpha$, then $\beth_{\alpha} = \sup_{\gamma < \alpha} \beth_{\gamma}$ is a supremum of $\abs{\alpha} < \mu$ many ordinals below $\mu$, so it is below $\mu$, as $\mu$ is regular.
\end{proof}

With these, all the axioms hold at once.

\begin{boxcorollary}
    If $M$ is strongly inaccessible, then $\sigma^{V_M}$ for all axioms $\sigma$ of $\ZFC$.
\end{boxcorollary}

\section{Relativization vs Satisfaction} \label{Ch3:Sec:Satisfaction}

Can we formalize into Set Theory the notion of ``$M \models \vp\of{\bar{x}}$''?

Pick ``symbols'' and formalize the following using $\ZFwoP$:
\begin{itemize}
    \item syntax of first-order logic
    \item definition of satisfaction
    \item proof theory
    \item completeness theorem
\end{itemize}

The way we define what ``formulae'' are internally is using Gödel numbers. Ie, for each formula $\vp$ (\textit{external}), there exists some set $\ulcorner \vp \urcorner$ that is analogous to it. We still write $M \models \vp\of{\bar{x}}$.

It is possible to show that our analogies translate both ways, ie, $\vp\of{\bar{x}}^M$ iff ``$M \models \vp\of{\bar{x}}$'', simply by structural induction on $\vp$.

% Transcribed from the board photographs supplied with the post-lecture pass of 8 Oct 2026; these boards were not typed up during the lecture.

Note, about $M \models \vp\of{\bar{x}}$:
\begin{enumerate}
    \item $\vp$ is not a set. There's a set $\ulcorner \vp \urcorner$ that is analogous to $\vp$. Nevertheless, we just write $M \models \vp\of{\bar{x}}$.
    \item This requires $M \ne \emptyset$ and $\bar{x} \in M^{n + 1}$, where $v_0, \ldots, v_n$ are the free variables of $\vp$. Contrast with $\vp^M\of{\bar{x}}$.
    \item $\vp^M\of{\bar{x}}$ was defined for each $\vp$ and each class $M$. $M \models \vp\of{\bar{x}}$ is defined only for sets $M$:
    \begin{align*}
        \models \; \ssq \; \setst{\parenth{M, c, \bar{x}}}{M \ne \emptyset, \, c \text{ a Gödel number}, \, \bar{x} \in M^{<\omega}}
    \end{align*}
\end{enumerate}

These agree wherever both make sense.

\begin{boxlemma}[Fact (Scheme)]
    Consider any formula $\vp\of{v_0, \ldots, v_n}$. For each $M \ne \emptyset$ and $\bar{x} \in M^{n + 1}$,
    \begin{align*}
        \vp\of{\bar{x}}^M \text{ iff } ``M \models \vp\of{\bar{x}}"
    \end{align*}
\end{boxlemma}

For now, work in $\ZF$. Given a formula $\vp(u)$ and a set $x$, say $\vp(u)$ \textbf{defines} $x$ iff
\begin{align*}
    \forall u \brac{u = x \lr \vp(u)} \tag{$\star$}
\end{align*}
Observe: there is no formula $\delta\of{r, x}$ so that, for every formula $\vp(u)$,
\begin{align*}
    \forall x \brac{\delta\of{\ulcorner \vp(u) \urcorner, x} \lr \underbrace{\vp(u) \text{ defines } x}_{(\star)}}
\end{align*}
Otherwise, the least undefinable ordinal would exist and be definable (because there are only countably many Gödel numbers).

\section{The Sets $H_{\kappa}$}

The levels $V_{\alpha}$ are cut out by rank. We now cut out sets by size instead---hereditarily, so that a set is small only if everything below it is. We will work in $\ZFC$ because we need choice to be able to work comfortably with cardinalities.

\subsection{Hereditary Cardinality}

We measure a set by the size of its transitive closure.

\begin{boxdefinition}[$H_{\kappa}$]
    For each cardinal $\kappa \geq \omega$, define the class
    \begin{align*}
        H_{\kappa} = \setst{z}{\abs{\trcl{\set{z}}} < \kappa}
    \end{align*}
\end{boxdefinition}

Note that
\begin{align*}
    \trcl{\set{z}} = \set{z} \cup \trcl{z}
\end{align*}
so
\begin{align*}
    \abs{\trcl{\set{z}}} = 1 \+ \abs{\trcl{z}}
\end{align*}

We can also unwind $\trcl{z}$ one level further, through the members of $z$.

% [CORRECTED] was: $\bigcup_{x \in z} \trcl{\set{z}}$ --- true, but only because the index $x$ never appears, so it says nothing; the uses below need $\trcl{\set{x}}$
\begin{boxproposition}
    For all sets $z$,
    \begin{align*}
        \trcl{\set{z}} = \set{z} \cup \bigcup_{x \in z} \trcl{\set{x}}
    \end{align*}
\end{boxproposition}

Nothing is assumed about $z$ here: it need not be transitive, let alone an ordinal.

\begin{boxlemma}
    Each $H_{\kappa}$ is a transitive class and, if $\kappa < \lambda$, then $H_{\kappa} \subsetneq H_{\lambda}$.
\end{boxlemma}
\begin{proof}
    \hfill
    \begin{description}
        \item[\underline{Transitivity.}] Let $z \in H_{\kappa}$ and $y \in z$. Then $y \in \trcl{\set{z}}$, so $\trcl{\set{z}}$ is a transitive set with $y$ as a member, and so it contains the smallest such set: $\trcl{\set{y}} \ssq \trcl{\set{z}}$. Hence $\abs{\trcl{\set{y}}} \leq \abs{\trcl{\set{z}}} < \kappa$, and $y \in H_{\kappa}$.

        \item[\underline{Monotonicity.}] If $\kappa < \lambda$, then $\abs{\trcl{\set{z}}} < \kappa$ implies $\abs{\trcl{\set{z}}} < \lambda$, so $H_{\kappa} \ssq H_{\lambda}$. The inclusion is strict because $\kappa \in H_{\lambda} \setminus H_{\kappa}$: $\trcl{\set{\kappa}} = \kappa + 1$ (it is transitive and has $\kappa$ as a member, and any transitive set with $\kappa$ as a member contains $\kappa$ and all its members), and $\abs{\kappa + 1} = \kappa$, as $\kappa$ is infinite, which is below $\lambda$ but not below $\kappa$.
    \end{description}
\end{proof}

In particular, the members of $H_{\kappa}$ have small rank.

\begin{boxlemma}
    For all $\kappa$, $H_{\kappa} \subseteq V_{\kappa}$.
\end{boxlemma}
\begin{proof}
    Let $z \in H_{\kappa}$. We will show that $z \in V_{\kappa}$. Let
    \begin{align*}
        \sigma = \setst{\rank{y}}{y \in \trcl{\set{z}}}
    \end{align*}
    Then $\sigma$ is a set of ordinals and $\rank{z} = \max(\sigma)$.

    We claim that $\sigma$ is a transitive set. Indeed, if not, then we would have some $\alpha < \beta$ such that $\beta = \min\of{\sigma \setminus \alpha}$. Pick $y \in \trcl{\set{z}}$ such that $\rank{y} = \beta$. Then, $\forall x \in y$, since $x \in \trcl{\set{y}}$ and $\rank{x} < \rank{y}$, we conclude that $\rank{x} < \alpha$. Then, by the rank formula,
    \begin{align*}
        \beta = \rank{y} = \sup_{x \in y}\of{\rank{x} + 1} \leq \alpha
    \end{align*}
    This is a contradiction! So $\sigma$ is transitive as desired.

    This then tells us that $\sigma = \rank{z} + 1$: a transitive set of ordinals is an ordinal, and an ordinal whose largest element is $\rank{z}$ is $\rank{z} + 1$.
    Then, we get
    \begin{align*}
        \abs{\sigma} \leq \sigma < \abs{\sigma}^+ \leq \abs{\trcl{\set{z}}}^+ \leq \kappa
    \end{align*}
    so $\rank{z} < \kappa$.
\end{proof}

So $H_{\kappa}$ is cut out of a set.

\begin{boxcorollary}
    Each $H_{\kappa}$ is a set.
\end{boxcorollary}
\begin{proof}
    Just apply comprehension.
\end{proof}

The first two of these get names of their own.

\begin{boxdefinition}
    Define $\HF := H_{\omega}$ and $\HC := H_{\omega_1}$
\end{boxdefinition}

We will show that $\HC$ is a model of $\ZFwoP$. In fact, so is $H_{\kappa^+}$ for each cardinal $\kappa$.

\begin{boxlemma}
    $\HF = V_{\omega}$
\end{boxlemma}
\begin{proof}
    We just saw that $\HF \ssq V_{\omega}$. Given some $z \in V_{\omega}$, pick $n < \omega$ such that $z \in V_n$. Then
    \begin{align*}
        \abs{\trcl{\set{z}}} \leq \abs{V_n} < \omega
    \end{align*}
    so $z \in \HF$.
\end{proof}

The ordinals in $H_{\kappa}$ are the ones we would expect.

\begin{boxlemma}
    $\OR \cap H_{\kappa} = \kappa$
\end{boxlemma}
\begin{proof}
    If $\alpha \in \OR$, then $\trcl{\set{\alpha}} = \alpha + 1$, because $\alpha + 1 = \alpha \cup \set{\alpha}$ is transitive and has $\alpha$ as a member, and any transitive set with $\alpha$ as a member contains $\alpha$ and all its members. So $\alpha \in H_{\kappa}$ iff $\abs{\alpha + 1} < \kappa$.

    If $\alpha < \kappa$, then $\alpha + 1 < \kappa$ too, because $\kappa$ is an infinite cardinal and so a limit ordinal; so $\abs{\alpha + 1} \leq \alpha + 1 < \kappa$. If $\alpha \geq \kappa$, then $\abs{\alpha + 1} \geq \abs{\kappa} = \kappa$. So the ordinals in $H_{\kappa}$ are exactly those below $\kappa$.
\end{proof}

So $H_{\kappa}$ and $V_{\alpha}$ can only coincide when the indices do.

\begin{boxlemma}
    If $H_{\kappa} = V_{\alpha}$, then $\kappa = \alpha$. 
\end{boxlemma}

Here is a useful lemma.

\begin{boxlemma}
    If $\kappa$ is regular, then
    \begin{align*}
        H_{\kappa} = \setst{z}{\forall x \in \trcl{\set{z}}, \, \abs{x} < \kappa}
    \end{align*}
\end{boxlemma}
\begin{proof}
    Call the RHS $I$. We show $H_{\kappa} = I$.

    \begin{description}
        \item[\underline{$H_{\kappa} \ssq I$.}] Let $z \in H_{\kappa}$ and $x \in \trcl{\set{z}}$. Then $x \ssq \trcl{\set{z}}$ and $\abs{x} \le \abs{\trcl{\set{z}}} < \kappa$.

        \item[\underline{$H_{\kappa} \supseteq I$.}] We apply $\in$-induction. Ie, we assume that for $x \in z$, since $x \in I$ (by transitivity of $I$), $x \in H_{\kappa}$, and therefore, $\abs{\trcl{\set{x}}} < \kappa$.

        Performing the induction, we calculate
        \begin{align*}
            \abs{\trcl{\set{z}}} \le 1 \+ \bigoplus_{x \in z} \underbrace{\abs{\trcl{\set{x}}}}_{< \kappa}
        \end{align*}
        As $z \in I$, $\abs{z} < \kappa$. So we are summing $< \kappa$ many cardinals that are all $< \kappa$, resulting in something $< \kappa$.
    \end{description}
\end{proof}

\subsection{The Axioms in $H_{\kappa}$}

So now let's ask ourselves what Axioms of ZFC hold in these $H_{\kappa}$.

\begin{boxlemma}
    The following ZFC axioms hold in all $H_{\kappa}$:
    \begin{enumerate}
        \item Extensionality
        \item Comprehension
        \item Pairing
        \item Union
        \item Choice
        \item Foundation \textit{(if we assume foundation in $V$)}
    \end{enumerate}
\end{boxlemma}
\begin{proof} \hfill
    \begin{enumerate}
        \item $H_{\kappa}$ is transitive.
        \item If $A \in H_{\kappa}$ and $B \ssq A$, then $B \in H_{\kappa}$.
        % [CORRECTED] items 3 and 4 were both: ``Essentially just because $H_{\kappa}$ is a set.'' --- that a class is a set gives no member of it containing $x$ and $y$ (eg, $2 = \set{0, 1}$ is a transitive set in which Pairing fails); what is needed is closure of $H_{\kappa}$ under pairs and unions
        \item If $x, y \in H_{\kappa}$, then $\set{x, y} \in H_{\kappa}$, since $\trcl{\set{\set{x, y}}} = \set{\set{x, y}} \cup \trcl{\set{x}} \cup \trcl{\set{y}}$ has cardinality below $\kappa$, as $\kappa$ is infinite.
        \item If $A \in H_{\kappa}$, then $\bigcup A \in H_{\kappa}$, since $\bigcup A \ssq \trcl{A}$, and so $\trcl{\set{\bigcup A}} \ssq \set{\bigcup A} \cup \trcl{A}$.
        \item Let $A \in H_{\kappa}$. Pick a well-ordering $R$ of $A$ (assuming AC in $V$). Then, as $R \ssq A \times A$, $R \in H_{\kappa}$. Use downwards $\Pi_1$-absoluteness to see that
        \begin{align*}
            \brac{(A, R) \text{ is a well-ordering}}^{H_{\kappa}}
        \end{align*}
        % the lecture did not argue item 6, Foundation
        \item \sorry
    \end{enumerate}
\end{proof}

Infinity needs $\kappa$ to be large enough to contain $\omega$.

\begin{boxlemma}
    Infinity holds in $H_{\kappa}$ iff $\kappa > \omega$.
\end{boxlemma}
\begin{proof}
    Follows from the fact that $\kappa > \omega \lr \omega \in H_{\kappa}$.
\end{proof}

Replacement is where the regularity of $\kappa$ comes in.

\begin{boxlemma} \label{Ch3:Lemma:H_Kappa_Replacement}
    If $\kappa = \cf(\kappa)$ (ie, $\kappa$ is regular), then each axiom $\sigma$ of the Replacement scheme holds relativized to $H_{\kappa}$, ie, $\sigma^{H_{\kappa}}$.
\end{boxlemma}
\begin{proof}
    If $A \in H_{\kappa}$ and $f : A \to H_{\kappa}$, then, letting $B = \range\of{f}$, we have
    \begin{align*}
        \abs{\trcl{\set{B}}}
        &= \abs{\set{B} \cup \bigcup_{x \in A} \trcl{\set{f(x)}}} \\
        &\leq 1 \+ \underbrace{\sum_{x \in A}}_{\text{fewer than } \kappa \text{ terms}} \underbrace{\abs{\trcl{\set{f(x)}}}}_{\text{each } < \kappa} \\
        &< \kappa
    \end{align*}
    There are strictly fewer than $\kappa$ terms because $A \ssq \trcl{\set{A}}$, and a sum of fewer than $\kappa$ cardinals, each below $\kappa$, is below $\kappa$ because $\kappa$ is regular. So $B \in H_{\kappa}$.

    This is exactly what Replacement in $H_{\kappa}$ needs. If $A, \bar{z} \in H_{\kappa}$ and each $x \in A$ has a unique $y \in H_{\kappa}$ with $\vp^{H_{\kappa}}\of{x, y, \bar{z}}$, let $f(x)$ be that $y$; $f$ is a set by Replacement in $V$, as $A \ssq H_{\kappa}$, and $B = \range\of{f} \in H_{\kappa}$ is the witness the axiom asks for.
\end{proof}

Power Set, on the other hand, fails at every successor cardinal.

\begin{boxlemma}
    If $\lambda = \kappa^+$, then
    \begin{align*}
        \parenth{\neg \text{(Power Set)}}^{H_{\lambda}}
    \end{align*}
\end{boxlemma}
\begin{proof}
    $\Powset(\kappa) \ssq H_{\lambda}$ but $\Powset(\kappa) \notin H_{\lambda}$.
\end{proof}

In fact, we can say exactly when it holds.

\begin{boxlemma} \label{Ch3:Lemma:H_Kappa_Power_Set}
    TFAE:
    \begin{enumerate}[label = (\arabic*)]
        \item $(\text{Power Set})^{H_{\lambda}}$
        \item $\lambda$ is a strong limit
    \end{enumerate}
\end{boxlemma}
\begin{proof} \hfill
    \begin{description}
        \item[\underline{$(2) \to (1)$.}] Assume $(2)$. Let $A \in H_{\lambda}$. Then
        \begin{align*}
            \abs{\trcl{\set{\Powset(A)}}}
            = \abs{\set{\Powset(A)} \cup \Powset(A) \cup \trcl{A}}
            \leq 1 \+ 2^{\abs{A}} \+ \abs{\trcl{A}}
            < \lambda
        \end{align*}

        \item[\underline{$\neg(2) \to \neg(1)$.}] We show, effectively, that if $\kappa < \lambda$ is a cardinal with $2^{\kappa} \geq \lambda$, then $\kappa \in H_{\lambda}$ and $\Powset(\kappa) \ssq H_{\lambda}$ but $\Powset(\kappa) \notin H_{\lambda}$, so Power Set fails in $H_{\lambda}$.
    \end{description}
\end{proof}

Putting all this together, we get the following.

\begin{boxcorollary} \hfill
    \begin{enumerate}
        \item $H_{\omega} = V_{\omega}$ models $\ZFCwoI$
        \item $H_{\omega_1} = H_{\omega_1}$ and all other $H_{\omega_2}, H_{\omega_3}, \ldots$ model $\ZFCwoP$
    \end{enumerate}
\end{boxcorollary}

The first cardinal not covered by the corollary is $\aleph_{\omega}$, which is singular, and it shows that the regularity hypothesis in \Cref{Ch3:Lemma:H_Kappa_Replacement} cannot be dropped.

\begin{boxexample}[$H_{\aleph_{\omega}}$]
    Everything but Replacement and Power Set holds in $H_{\aleph_{\omega}}$: Extensionality, Comprehension, Pairing, Union, Choice and Foundation hold in every $H_{\kappa}$, and Infinity holds because $\aleph_{\omega} > \omega$.

    Replacement fails. There's a formula $\vp\of{u, v}$ so that, for all $n \in \omega$ and $\kappa \in H_{\aleph_{\omega}}$,
    \begin{align*}
        \vp\of{n, \kappa}^{H_{\aleph_{\omega}}} \lr \kappa = \aleph_n
    \end{align*}
    Namely, $\vp\of{n, \kappa}$ says that there is a function $g$ with domain $n + 1$ such that $g(0) = \omega$, each $g(i + 1)$ is the least cardinal above $g(i)$, and $g(n) = \kappa$. This relativizes correctly because the witness $\grpres{\aleph_i}{i \leq n}$ lies in $H_{\aleph_{\omega}}$, and because an ordinal below $\aleph_{\omega}$ is a cardinal in $H_{\aleph_{\omega}}$ iff it is one in $V$, as every bijection between two ordinals below $\aleph_{\omega}$ lies in $H_{\aleph_{\omega}}$. So $\vp$ defines $f : n \mapsto \aleph_n$ over $H_{\aleph_{\omega}}$, on the domain $\omega \in H_{\aleph_{\omega}}$, but $f \notin H_{\aleph_{\omega}}$.
    \begin{figure}[H]
        \centering
        \begin{tikzpicture}[dot/.style = {circle, fill, inner sep = 1.5pt}, every node/.style = {font = \small}]
            \drawcone{2}{4.5}{$H_{\aleph_{\omega}}$}
            \node[dot, label = {left:$\omega$}] (w) at (0.1, 0.7) {};
            \node[dot, label = {left:$\aleph_n$}] (a) at (-0.5, 3.2) {};
            \draw[->, thick] (w) to[bend right = 60] node[right] {$f : n \mapsto \aleph_n$} (a);
        \end{tikzpicture}
    \end{figure}
    Indeed, any $B$ with every $\aleph_n \in B$ has $\aleph_{\omega} = \bigcup_{n < \omega} \aleph_n \ssq \trcl{\set{B}}$, so $\abs{\trcl{\set{B}}} \geq \aleph_{\omega}$ and $B \notin H_{\aleph_{\omega}}$. So no set in $H_{\aleph_{\omega}}$ collects the values of $f$.

    Whether Power Set holds in $H_{\aleph_{\omega}}$ is, by \Cref{Ch3:Lemma:H_Kappa_Power_Set}, the question of whether $\aleph_{\omega}$ is a strong limit, which $\ZFC$ does not settle.
\end{boxexample}

\subsection{$\Sigma_1$-Correctness}

The most useful property of the $H_{\kappa}$ is that they are right about every $\Sigma_1$ statement, and proving it uses the satisfaction relation of \Cref{Ch3:Sec:Satisfaction}.

\begin{boxtheorem} \label{Ch3:Thm:H_Kappa_Sigma1}
    If $\kappa$ is an uncountable cardinal, then $H_{\kappa} \prec_{\Sigma_1} V$.
\end{boxtheorem}

\begin{remark}
    This means that for each $\Sigma_1$ formula $\vp\of{\bar{u}}$ and each matching $\bar{x} \in H_{\kappa}$, $\vp\of{\bar{x}}^{H_{\kappa}} \lr \vp\of{\bar{x}}^V$. In other words, it means $\Sigma_1$-formulae are absolute for $H_{\kappa}$. We also say, ``$H_{\kappa}$ is $\Sigma_1$-correct in $V$.''
\end{remark}

Note that $H_{\kappa} \prec_{\Sigma_1} V$ means $\Sigma_1$ formulas are absolute between $H_{\kappa}$ and $V$. We borrow model theoretic notation for being an elementary substructure, but note that this is the model theory we know in the ``outside world'', not the model theory we have inside of $V$.

\begin{proof}
    We already know upwards absoluteness, so the only remaining question is downwards absoluteness.

    Suppose $\vp(u)$ is the formula $\exists v\parenth{\psi\of{u, v}}$, with $\psi$ $\Delta_0$. Suppose $x \in H_{\kappa}$ is so that $\vp\of{x}$.

    We need to show that
    \begin{align*}
        \parenth{\exists v\parenth{\psi\of{x, v}}}^{H_{\kappa}}
    \end{align*}

    Pick any $y$ such that $\psi\of{x, y}$. Let $N = \trcl{\set{x, y}}$. Then $x, y \in N$ and $N$ is transitive. By $\Delta_0$ absoluteness, we know $\psi(x, y)^N$. Can write ``$N \models \psi(x, y)$''.

    % [FILLED] the proof from here on is supplied; the lecture broke off at the application of Downwards Löwenheim-Skolem
    Apply Downwards Löwenheim-Skolem to find $Z \prec N$ with $\trcl{\set{x}} \ssq Z$ and $\abs{Z} \leq \abs{\trcl{\set{x}}} + \aleph_0$. As $x \in H_{\kappa}$ and $\kappa$ is uncountable, $\abs{Z} < \kappa$.

    $N$ is transitive, so Extensionality holds in $N$, and hence in $Z \prec N$; so $\in$ is extensional on $Z$. It is also well-founded and set-like there, so by the Mostowski Collapse Theorem Scheme (\Cref{Ch2:Thm:Mostowski}) there is a transitive set $M$ and an isomorphism $\pi : \parenth{Z, \in} \to \parenth{M, \in}$, given by $\pi(z) = \setst{\pi(w)}{w \in z \cap Z}$. Two things about $\pi$:
    \begin{itemize}
        \item $\pi(w) = w$ for every $w \in \trcl{\set{x}}$. Indeed, by $\in$-induction: such a $w$ is a subset of the transitive set $\trcl{\set{x}} \ssq Z$, so $w \cap Z = w$, and $\pi(w) = \setst{\pi(v)}{v \in w} = w$. In particular, $\pi(x) = x$.
        \item $\abs{M} = \abs{Z} < \kappa$.
    \end{itemize}

    Now $N \models \exists v \, \psi\of{x, v}$, so, as $Z \prec N$ and $x \in Z$, there is $y' \in Z$ with $Z \models \psi\of{x, y'}$. Applying the isomorphism, $M \models \psi\of{x, \pi(y')}$, ie, $\psi\of{x, \pi(y')}^M$. As $M$ is transitive and $\psi$ is $\Delta_0$, $\Delta_0$-absoluteness gives $\psi\of{x, \pi(y')}$.

    Finally, $\pi(y') \in M$ and $M$ is transitive, so $\trcl{\set{\pi(y')}} \ssq M$ and $\abs{\trcl{\set{\pi(y')}}} \leq \abs{M} < \kappa$. So $\pi(y') \in H_{\kappa}$, and as $H_{\kappa}$ is transitive, $\Delta_0$-absoluteness gives $\psi\of{x, \pi(y')}^{H_{\kappa}}$. So $\parenth{\exists v \, \psi\of{x, v}}^{H_{\kappa}}$.
\end{proof}

This has a consequence for the levels of the $V$-hierarchy, through the $\beth$-fixed points.

Recall that $\setst{\kappa}{\kappa = \beth_{\kappa}}$ is a club class of cardinals. It is closed because the $\beth$ operation is ``continuous'': if $\lambda$ is a limit of $\beth$-fixed points, then, as $\alpha \mapsto \beth_{\alpha}$ is increasing, $\beth_{\lambda} = \sup_{\alpha < \lambda} \beth_{\alpha}$ is also the supremum of $\beth_{\kappa} = \kappa$ over the fixed points $\kappa < \lambda$, which is $\lambda$. So a sup of $\beth$-fixed points is a $\beth$-fixed point. It is unbounded because, given any $\alpha$, the sequence $\kappa_0 = \alpha$, $\kappa_{n + 1} = \beth_{\kappa_n}$ is non-decreasing (as $\beth_{\gamma} \geq \gamma$) and its supremum $\lambda \geq \alpha$ is a fixed point: either the sequence is eventually constant, at a fixed point, or $\lambda$ is a limit and $\beth_{\lambda} = \sup_{n} \beth_{\kappa_n} = \sup_n \kappa_{n + 1} = \lambda$.

We also saw that if $\mu$ is a strongly inaccessible cardinal, then $\mu = \beth_{\mu}$ (\Cref{Ch3:Prop:Inaccessible_Beth}).

We have seen that $H_{\kappa} \prec_{\Sigma_1} V$ for all uncountable cardinals $\kappa$. It follows that for such $\kappa$, $\kappa = \beth_{\kappa}$ iff $H_{\kappa} = V_{\kappa}$, which in turn implies that $V_{\kappa} \prec_{\Sigma_1} V$.

Observe that $V_{\kappa} \prec_{\Sigma_1} V$ is \underline{not} a formula in the language of set theory. But $\kappa = \beth_{\kappa}$ and $H_{\kappa} = V_{\kappa}$ \textit{are}. By comprehension, put
\begin{align*}
    C_1 = \setst{\kappa}{\kappa = \beth_{\kappa}} = \setst{\kappa}{\kappa \text{ is uncountable and } H_{\kappa} = V_{\kappa}}
\end{align*}
Then $C_1$ is a class. Writing $C_1 = \setst{\kappa}{V_{\kappa} \prec_{\Sigma_1} V}$ is problematic.

Note that $C_1$ is club in $\OR$. Indeed, we proved that the set of $\beth$-fixed points contains a club class of ordinals. Define a class function $f : \OR \to \OR$ by
\begin{itemize}
    \item $f(0) = \beth_0$
    \item $f(\alpha + 1) = \beth_{f(\alpha) + 1}$
    \item $f(\beta) = \sup_{\alpha < \beta} f(\alpha)$
\end{itemize}
% [CORRECTED] was: ``for every limit $\beta \in \lim\of{\range\of{f}}$, then $\beta = f(\beta) = \beth_{\beta}$'' --- $f$ is strictly increasing, so $f(\beta) > \beta$ in general (eg, $\beta = f(\omega)$); what holds is that the values of $f$ at limits, which are exactly the limit points of its range, are $\beth$-fixed points
The range of $f$ is a cub and, for every limit ordinal $\beta$, $f(\beta) = \beth_{f(\beta)}$; that is, every element of $\lim\of{\range\of{f}}$\footnote{This is the class of limit points of $\range(f)$, and as we know, the class of limit points of a club is a club} is a $\beth$-fixed point. So $C_1$ contains a club. So $C_1$ is unbounded. Moreover, $C_1$ is closed. So it's a club in $\OR$.

\section{Reflection}

Below an inaccessible, we can ask for more than $\Sigma_1$-correctness: since $V_{\mu}$ is a set, full elementarity makes sense.

\subsection{Reflection Below an Inaccessible}

Here is the precise statement.

\begin{boxtheorem}
    Suppose $\mu$ is a strongly inaccessible cardinal. Let $C = \setst{\kappa < \mu}{V_{\kappa} \prec V_{\mu}}$. Then $C$ is a club subset of $\mu$.
\end{boxtheorem}

\begin{remark}
    $V_{\kappa} \prec V_{\mu}$ \textbf{is} a formula in the language of set theory because both $V_{\kappa}$ and $V_{\mu}$ are \underline{sets}. It would be a different story if they were classes.

    ``This is model theory. Model theory is math. And math is embedded in set theory. So this makes sense!''
\end{remark}

\begin{proof}
    Clearly $C$ is closed in $\mu$. If $C \cap \lambda \ne \emptyset$ and $\lambda = \sup(C \cap \lambda)$, then we need $\lambda \in C$ (meaning of closed).

    For such a $\lambda$, we have $V_{\lambda} = \bigcup_{\kappa < \lambda} V_{\kappa} = \bigcup_{\kappa \in C \cap \lambda} V_{\kappa}$.

    \begin{itemize}
        \item It follows, for $\kappa < \kappa' < \mu$ with $\kappa, \kappa' \in C$, that $V_{\kappa} \prec V_{\kappa'}$.

        \item Then, $\grpres{V_{\kappa}}{\kappa \in C \cap \lambda}$ is a chain of elementary substructures of $V_{\mu}$, so its union is also an elementary substructure of $V_{\mu}$.

        \item $V_{\lambda} \prec V_{\mu}$, so $\lambda \in C$.

        \item It also follows that $V_{\kappa} \prec V_{\lambda}$ for $\kappa \in C \cap \lambda$.
    \end{itemize}

    Now let $\kappa_0 < \mu$ be an arbitrary uncountable cardinal. By Downwards Löwenheim-Skolem, we have $X_0 \prec V_{\mu}$ with $V_{\kappa_0} \cup \set{\kappa_0} \ssq X_0$ and $\abs{X_0} = \abs{V_{\kappa_0} \cup \set{\kappa_0}} = \beth_{\kappa_0} < \beth_{\mu} = \mu$.

    % IDEA: Take X_0, fill in things of smaller rank, then take X_1, fill in evrything of smaller rank, and so on. Each X_i is obtained by applying DLS to $V_{\mu}$. We then get $V_{\kappa} \bigcup_{n < \omega} V_{\kappa_n} = \bigcup_{n < \omega} X_n \prec V_{\kappa}$ where we konw the chain is elementary by the elementary chain theorem.
    \sorry
\end{proof}

\subsection{Hierarchies}

To state reflection in general, we give a name to the shape that $\grpres{V_{\alpha}}{\alpha \in \OR}$ and the $H_{\kappa}$ share.

\begin{boxdefinition}[Hierarchy]
    We say that a class function $\grpres{Z_{\alpha}}{\alpha \in C}$ is a \textbf{hierarchy} of sets iff
    \begin{itemize}
        \item each $Z_{\alpha}$ is a transitive set
        \item $C$ is a club class of ordinals
        \item for every $\alpha < \beta$ from $C$, $Z_{\alpha} \subset Z_{\beta}$
        \item for every $\beta$ a limit point\footnote{as $C$ is club, such $\beta$s do also lie in $C$} of $C$,
        \begin{align*}
            Z_{\beta} = \bigcup_{\alpha \in C \cap \beta} Z_{\alpha}
        \end{align*}
    \end{itemize}
    We usually write
    \begin{align*}
        Z = \bigcup_{\alpha \in C} Z_{\alpha}
        = \setst{x}{\exists \alpha \in C \st x \in Z_{\alpha}}
    \end{align*}
    (unless, of course, we already have a name for our hierarchy).
\end{boxdefinition}

We have already seen two of these.

\begin{boxexample}
    The following are all examples of hierarchies we've seen:
    \begin{align*}
        \grpres{V_{\alpha}}{\alpha \in \OR}
        \qquad , \qquad
        \grpres{H_{\kappa}}{\kappa \text{ an infinite cardinal}}
    \end{align*}
\end{boxexample}

Indeed, there will be examples where $Z \ne V$.

\subsection{The Reflection Theorem}

Every hierarchy reflects any finitely many formulae down to club many of its levels.

\begin{boxtheorem}[Reflection theorem scheme, proved in $\ZF$]
    Suppose $\varphi_i$ for $i < \ell < \omega$ are formulae in the language of set theory. Suppose that $\grpres{Z_{\alpha}}{\alpha \in C}$ is a hierarchy. Then there's a club class $D \ssq C$ such that for every $\alpha \in D$ and $i < \ell$, $\vp_i$ is absolute between $Z_{\alpha}$ and $Z$.
\end{boxtheorem}
\begin{proof}
    WLOG, assume that for all $j < \ell$ and proper subformulae $\psi$ of $\varphi_j$, there is some $i < j$ such that $\psi$ is $\vp_{i}$. Ie, we are setting up an induction on the complexity of $\vp$ by setting up an induction on ordinals.

    Let's define a class sequence of ordinals $\grpres{\alpha_{\eta}}{\eta \in \OR}$ as follows.
    \begin{itemize}
        % [CORRECTED] was: $\alpha_0 = 0$ --- $Z_{\alpha}$ is only defined for $\alpha \in C$, and $0$ need not be in $C$ (it is not for $\grpres{H_{\kappa}}{\kappa \text{ an infinite cardinal}}$), while the successor step uses $Z_{\alpha_0}$; the board also says to maintain $\alpha_{\eta} \in C$ for all $\eta$
        \item \underline{Base:} Define $\alpha_0 = \min(C)$.
        \item \underline{Successor:} Given $\alpha_{\eta}$, define $\alpha_{\eta + 1}$ to be the least $\beta \in C \setminus \parenth{\alpha_{\eta} + 1}$ so that for all $j < \ell$, if $\vp_{j}\of{\bar{u}}$ is of the form $\exists v, \vp_i\of{\bar{u}, v}$, then for all $\bar{x}$ from $Z_{\alpha_{\eta}}$ that match $\bar{u}$, if there is $y \in Z$ such that $\vp_i\of{\bar{x}, y}^Z$, then there is $y \in Z_{\beta}$ such that $\vp_i\of{\bar{x}, y}^Z$.
        \item \underline{Limit case:} If $\theta$ is a limit ordinal, define
        \begin{align*}
            \alpha_{\theta} = \sup_{\eta < \theta} \of{\alpha_{\eta}} \in \lim(C) \subset C
        \end{align*}
        (because $C$ is club).
    \end{itemize}

    The successor step is where the work happens. A tuple $\bar{x}$ from $Z_{\alpha_{\eta}}$ may have a witness $y'$ only higher up, in some $Z_{\beta'}$, and $\alpha_{\eta + 1}$ is chosen high enough to catch a witness for every such $\bar{x}$ at once:
    \begin{figure}[H]
        \centering
        \begin{tikzpicture}[dot/.style = {circle, fill, inner sep = 1.5pt}, every node/.style = {font = \small}]
            \drawcone{2.2}{5}{$Z$}
            \drawconelevel{2.2}{5}{1.6}{$Z_{\alpha_{\eta}}$}
            \drawconelevel{2.2}{5}{4.2}{$Z_{\beta'}$}
            \node[dot, label = {right:$\bar{x}$}] at (-0.15, 1.0) {};
            \node[dot, label = {right:$y'$}] at (0.3, 3.6) {};
        \end{tikzpicture}
    \end{figure}
    Only existential formulae need this treatment. We may take every $\vp_j$ to be built from atomic formulae using $\neg$, $\land$ and $\exists$ alone, writing $\forall$ as $\neg \exists \neg$. Relativization leaves atomic formulae unchanged, so they mean the same in $Z_{\alpha}$ as in $Z$, and it commutes with $\neg$ and $\land$, so absoluteness passes through them. The one place it can fail is $\exists v$, where $Z$ may have a witness that $Z_{\alpha}$ lacks.

    The successor step makes sense. For each $j < \ell$ of the form above and each $\bar{x}$ from $Z_{\alpha_{\eta}}$ that has a witness $y \in Z$, there is a least $\gamma \in C$ with a witness in $Z_{\gamma}$, as $Z = \bigcup_{\gamma \in C} Z_{\gamma}$. There are only set many such pairs $\parenth{j, \bar{x}}$, as $Z_{\alpha_{\eta}}$ is a set, so these $\gamma$ have a supremum, and as $C$ is unbounded there is some $\beta \in C$ above it and above $\alpha_{\eta}$. Any such $\beta$ works, since $Z_{\gamma} \ssq Z_{\beta}$ for $\gamma \leq \beta$ in $C$.

    Then take $D$ to be $\setst{\alpha_{\theta}}{\theta \text{ is a limit ordinal}}$. The sequence $\grpres{\alpha_{\eta}}{\eta \in \OR}$ is strictly increasing and continuous at limits, so $D$ is a club class: it is unbounded since $\alpha_{\theta} \geq \theta$, and closed since a supremum of $\alpha_{\theta}$s over limit ordinals $\theta$ in some set $\Theta$ with no largest element is $\alpha_{\sup \Theta}$, and $\sup \Theta$ is a limit ordinal. And $D \ssq C$ by the limit case.

    Finally, fix $\alpha = \alpha_{\theta} \in D$; we show that each $\vp_i$ is absolute between $Z_{\alpha}$ and $Z$, by induction on $i$. The atomic and propositional cases are as above, so suppose $\vp_j\of{\bar{u}}$ is $\exists v \, \vp_i\of{\bar{u}, v}$ with $i < j$, and let $\bar{x}$ from $Z_{\alpha}$ match $\bar{u}$.
    \begin{description}
        \item[\underline{From $Z_{\alpha}$ to $Z$.}] If $y \in Z_{\alpha}$ and $\vp_i\of{\bar{x}, y}^{Z_{\alpha}}$, then $y \in Z$ and $\vp_i\of{\bar{x}, y}^Z$ by the induction hypothesis, so $\vp_j\of{\bar{x}}^Z$.

        \item[\underline{From $Z$ to $Z_{\alpha}$.}] Suppose $\vp_j\of{\bar{x}}^Z$. As $\alpha_{\theta}$ is a limit point of $C$, $Z_{\alpha_{\theta}} = \bigcup_{\gamma \in C \cap \alpha_{\theta}} Z_{\gamma} = \bigcup_{\eta < \theta} Z_{\alpha_{\eta}}$, and $\bar{x}$ is a finite tuple, so $\bar{x}$ is from $Z_{\alpha_{\eta}}$ for some $\eta < \theta$. By the choice of $\alpha_{\eta + 1}$, there is $y \in Z_{\alpha_{\eta + 1}} \ssq Z_{\alpha}$ with $\vp_i\of{\bar{x}, y}^Z$, and by the induction hypothesis $\vp_i\of{\bar{x}, y}^{Z_{\alpha}}$. So $\vp_j\of{\bar{x}}^{Z_{\alpha}}$.
    \end{description}
\end{proof}

Here's an interesting corollary.

Let $\sigma_0, \ldots, \sigma_{\ell}$ be finitely many axioms of $\ZFC$. By the reflection theorem scheme, there are club many $\alpha \in \OR$ so that $\forall  i \leq \ell$, $\sigma_i^{V_{\alpha}}$. Equivalently, $\forall i \leq \ell$, $V_{\alpha} \models \sigma_i$. Rephrasing:

\begin{boxcorollary}
    For every finite sub-theory $T$ of $\ZFC$,
    \begin{align*}
        \ZFC \vdash \text{``$T$ has a transitive model''}
    \end{align*}
    where the quotes (``read with a different tone of voice'')
\end{boxcorollary}

With Choice, we can make the model countable as well. Given $\alpha$ with
\begin{align*}
    V_{\alpha} \models \bigwedge_{i \leq \ell} \sigma_i
\end{align*}
we can apply Downwards Löwenheim-Skolem to get some countable $X \prec V_{\alpha}$, so that $X \models \bigwedge_{i \leq \ell} \sigma_i$ too. As $V_{\alpha}$ is transitive, it satisfies Extensionality, and hence so does $X$; so $\in$ is extensional on $X$, as well as well-founded and set-like. Then apply Mostowski (\Cref{Ch2:Thm:Mostowski}) to get $X \simeq M$ with $M$ transitive and countable, and $M \models \bigwedge_{i \leq \ell} \sigma_i$, because satisfaction is preserved by isomorphisms. So in fact $\ZFC$ proves that every finite sub-theory of $\ZFC$ has a countable transitive model.

% This file is the inbox for the lecture in progress. Type today's raw notes straight into
% it: it is \input from the chapter file, so the Overleaf preview builds them live. Once the
% material has been placed into the sections, /integrate empties the file again and leaves
% this header behind. Do not delete the file or its \input line.

\section{Absoluteness}

We will constantly be comparing what a formula says in a smaller class with what it says in a bigger one, so we give the possible outcomes names.

\begin{boxdefinition}[Absoluteness]
    Let $M \ssq N$ be classes and let $\vp\of{\bar{u}}$ be a formula. We say $\vp\of{\bar{u}}$ is
    \begin{itemize}
        \item \textbf{downwards absolute} from $N$ to $M$ iff for all $\bar{x} \in M^{<\omega}$ matching $\bar{u}$, $\vp^N\of{\bar{x}} \to \vp^M\of{\bar{x}}$.
        \item \textbf{upwards absolute} from $M$ to $N$ iff for all $\bar{x} \in M^{<\omega}$ matching $\bar{u}$, $\vp^M\of{\bar{x}} \to \vp^N\of{\bar{x}}$.
        \item \textbf{absolute} between $M$ and $N$ iff it is both, ie, iff for all $\bar{x} \in M^{<\omega}$ matching $\bar{u}$, $\vp^M\of{\bar{x}} \lr \vp^N\of{\bar{x}}$.
    \end{itemize}
    When $N = V$, we just say $\vp\of{\bar{u}}$ is downwards absolute, upwards absolute or absolute \textbf{for $M$}.
\end{boxdefinition}

\subsection{Absoluteness Theorems}

Every absoluteness theorem we prove will have an assumption that the smaller model (or sometimes, both models) are transitive. There are no absoluteness theorems in this course that do not mention transitivity, and it is a common mistake to try and apply them in cases where there isn't transitivity.

\begin{boxlemma}[$\Delta_0$-Absoluteness]
    Let $M$ and $N$ be classes\footnote{possibly proper classes, possibly sets}, with $M \ssq N$. Assume $M$ is transitive. Then, for every $\Delta_0$ formula $\varphi\of{\bar{u}}$ and $\bar{x} \in M^{<\omega}$ (matching $\bar{u}$ in length),
    \begin{align*}
        \varphi^M\of{\bar{x}} \lr \varphi^N\of{\bar{x}}
    \end{align*}
    That is, $\varphi\of{\bar{u}}$ is absolute between $M$ and $N$.
\end{boxlemma}
\begin{proof}
    We argue by induction on the complexity of $\varphi\of{\bar{u}}$.

    \begin{itemize}
        \item \underline{Atomic Case.} Note that
        \begin{align*}
            \parenth{u_0 = u_1}^M
            \qquad
            u_0 = u_1
            \qquad
            \parenth{u_0 = u_1}^N
        \end{align*}
        are all exactly the same. Similarly,
        \begin{align*}
            \parenth{u_0 \in u_1}^M
            \qquad
            u_0 \in u_1
            \qquad
            \parenth{u_0 \in u_1}^N
        \end{align*}
        are all exactly the same.

        \item \underline{Propositional Case.} If $\varphi$ is some propositional combination of $\varphi_0$ and $\varphi_1$, then the statement for $\vp$ is clear from the induction hypotheses for $\vp_0$ and $\vp_1$.

        \item \underline{Bounded Quantifiers.} Say $\vp\of{\bar{u}}$ is of the form $\exists v \in w \, \parenth{\psi\of{\bar{u}, v, w}}$. (This is enough, because bounded universal can be formulated as a negation of a bounded existential, and the propositional case covers negations.)

        Let $\bar{x} \in M^{<\omega}$ and $z \in M$. Then, the following are equivalent:
        \begin{itemize}
            \item $\brac{\exists v \in z, \, \psi\of{\bar{x}, v, z}}^M$
            \item There is $y \in z \cap M$ such that $\psi^M\of{\bar{x}, y, z}$
            \item There is $y \in z$ such that $\psi^M\of{\bar{x}, y, z}$ (because $M$ is transitive, so $z = z \cap M$)
            \item There is $y \in z \cap N$ such that $\psi^M\of{\bar{x}, y, z}$ (because $z \ssq M \ssq N$, so $z = z \cap N$)
            \item There is $y \in z \cap N$ such that $\psi^N\of{\bar{x}, y, z}$
            \item $\brac{\exists v \in z, \, \psi\of{\bar{x}, v, z}}^N$
        \end{itemize}
    \end{itemize}
\end{proof}

\begin{boxlemma}
    In the same situation as above, we have upward absoluteness from $M$ to $N$ for $\Sigma_1$ formulae and downward absoluteness from $N$ to $M$ for $\Pi_1$ formulae.
\end{boxlemma}

\begin{boxlemma}
    Let $T$ be a theory and $M \ssq N$ be models of $T$ (with $M$ transitive). Then,
    \begin{itemize}
        \item $\Sigma_1^T$ formulae are upwards-absolute
        \item $\Pi_1^T$ formulae are downwards-absolute
        \item $\Delta_1^T$ formulae are absolute
    \end{itemize}
\end{boxlemma}

\subsection{Absoluteness of ZFC Axioms}

For this discussion, fix a class $M$. Let's list out the axioms of $\ZFwoF$, relativised to $M$.

We start with Extensionality, which relativized to $M$ reads
\begin{align*}
    \brac{\forall A \forall B \brac{\forall x \parenth{x \in A \lr x \in B} \to A = B}}^M
\end{align*}
Or, equivalently,
\begin{align*}
    \brac{\forall A \forall B \brac{\parenth{A \ssq B \land B \ssq A} \to A = B}}^M
\end{align*}
where $A \ssq B$ is the $\Delta_0$ formula $\parenth{\forall x \in A}\parenth{x \in B}$. Indeed, all sub-formulae of the above are $\Delta_0$.

This is all we need, as long as $M$ is transitive.

\begin{boxlemma}[Extensionality]
    If $M$ is transitive, then $(\text{Extensionality})^M$.
\end{boxlemma}
\begin{proof}
    By $\Delta_0$-absoluteness, for $A, B \in M$ the formula $\parenth{A \ssq B \land B \ssq A} \to A = B$ says the same thing in $M$ as it does in $V$, where it is true. So it holds in $M$ for all $A, B \in M$, which is exactly $(\text{Extensionality})^M$.
\end{proof}

From now on, assume that \underline{$M$ is transitive}, so that $(\text{Extensionality})^M$.

Comprehension needs $M$ to contain enough subsets of its members.

\begin{boxlemma}[Comprehension]
    If $M$ is closed under subsets, ie, $A \in M$ and $B \ssq A$ imply $B \in M$, then $(\text{Comprehension})^M$.
\end{boxlemma}
\begin{proof} % [FILLED] the analysis of Comprehension^M, and the closure condition that secures it, are supplied here; the lecture broke off after the transitivity assumption
    For each formula $\vp\of{x, \bar{y}}$, the corresponding instance relativized to $M$ is
    \begin{align*}
        \forall A \in M \, \forall \bar{y} \in M \, \exists B \in M \, \forall x \in M \brac{x \in B \lr \parenth{x \in A \land \vp^M\of{x, \bar{y}}}}
    \end{align*}
    As $M$ is transitive, any such $B$ and the given $A$ are subsets of $M$, so the bracket pins $B$ down completely: it must be $\setst{x \in A}{\vp^M\of{x, \bar{y}}}$. So Comprehension relativizes to $M$ iff each of these sets lies in $M$, and in particular it does whenever $M$ is closed under subsets.
\end{proof}

From now on, assume $(\text{Comprehension})^M$ as well.

Pairing only needs $M$ to be closed under pairs.

\begin{boxlemma}[Pairing]
    If $\set{x, y} \in M$ for all $x, y \in M$, then $(\text{Pairing})^M$.
\end{boxlemma}
\begin{proof} % [FILLED] the lecture left Pairing blank
    Relativized to $M$, Pairing says
    \begin{align*}
        \forall x \in M \, \forall y \in M \, \exists A \in M \, \parenth{x \in A \land y \in A}
    \end{align*}
    If $\set{x, y} \in M$, then it is a witness. (Here we did not even need transitivity, because the formula under the quantifiers is atomic.)
\end{proof}

We will want to relativize classes as well as formulae, and it is convenient to have notation for it.

\begin{boxnotation}
    Say $M$ is a transitive class and $A$ another class. Say $A = \setst{x}{\vp(x)}$. We might write $A^M$ to mean $\setst{x \in M}{\vp^M(x)}$ even though it may depend on the choice of $\varphi$ with this notation.

    $A^M$ is what $M$ thinks $A$ is. When it agrees with $A \cap M$, ie, when $\vp(x)$ is absolute for $M$, we say $A$ is \textbf{absolute} for $M$.
\end{boxnotation}

This is where the digression is going: below, we ask whether $\OR^M = \OR \cap M$, and whether a class function $G$ defined by recursion satisfies $G^M = G \restriction \parenth{A \cap M}$.

% 21 Sep 2026

Let's see which axioms hold in the levels of the $V$-hierarchy.

\begin{boxproposition}
    For each axiom $\sigma$ of $\set{\text{Extensionality}, \text{Comprehension}, \text{Pairing}, \text{Unions}, \text{Power Sets}}$, for all limit ordinals $\alpha$, $\sigma^{V_{\alpha}}$.
\end{boxproposition}

\begin{boxproposition}
    For each axiom $\sigma$ of $\set{\text{Extensionality}, \text{Comprehension}, \text{Pairing}, \text{Unions}, \text{Power Sets}, \text{Replacement}}$, $\sigma^{V_{\omega}}$.
\end{boxproposition}

\begin{boxnexample}[$V_{\omega + \omega}$]
    We can show that Replacement does not relativise to $V_{\omega + \omega}$. Indeed, consider the class function $G : n \mapsto \omega + n$. It's definable in two ways:
    \begin{enumerate}
        % [CORRECTED] was: $\exists g \brac{\dom(g) = n + 1 \land g(0) = w \land g(i + 1) = g(i) \cup \set{g(i) = g(i) + 1 \land g(n) = S}}$ --- $i$ was never quantified, the brace closed around $g(n) = S$, $w$ stood for $\omega$, and nothing said $g$ is a function, so it did not express $S = \omega + n$
        \item $G(n) = S$ iff $\exists g \brac{g \text{ is a function} \land \dom{g} = n + 1 \land g(0) = \omega \land \forall i < n \, \parenth{g(i + 1) = g(i) \cup \set{g(i)}} \land g(n) = S}$
        % [FILLED] the second definition and the argument after it are supplied here; the lecture announced two definitions, gave one, and stopped
        \item $G(n) = S$ iff $\forall g \brac{\parenth{g \text{ is a function} \land \dom{g} = n + 1 \land g(0) = \omega \land \forall i < n \, \parenth{g(i + 1) = g(i) \cup \set{g(i)}}} \to g(n) = S}$
    \end{enumerate}
    Both say that $S$ is the last value of the function $g$ on $n + 1$ that starts at $\omega$ and takes a successor at each step, and in both, everything after the unbounded quantifier over $g$ is $\Delta_0$ in $g$, $n$, $S$ and the parameter $\omega$. So the first is $\Sigma_1$ and the second is $\Pi_1$.

    Now fix $n < \omega$. The unique $g$ in question is $\setst{\parenth{i, \omega + i}}{i \leq n}$, and each of its pairs lies in $V_{\omega + n + 3}$, so $g \in V_{\omega + n + 4} \ssq V_{\omega + \omega}$. Hence, for $S \in V_{\omega + \omega}$, the first formula relativized to $V_{\omega + \omega}$ holds iff $S = \omega + n$: the witness $g$ is available inside $V_{\omega + \omega}$, and the $\Delta_0$ part means the same thing there as in $V$ by $\Delta_0$-absoluteness. (The second formula relativizes correctly too, by downward absoluteness of $\Pi_1$ formulae, so it does not matter which definition we use.)

    So the hypothesis of the instance of Replacement for this formula, with $A = \omega \in V_{\omega + \omega}$ and parameter $\omega$, holds in $V_{\omega + \omega}$. If the conclusion did too, we would have $B \in V_{\omega + \omega}$ with $\omega + n \in B$ for every $n < \omega$. But $B \in V_{\omega + k + 1}$ for some $k < \omega$, so $B \ssq V_{\omega + k}$, and then $\omega + k \in V_{\omega + k}$, ie, $\rank{\omega + k} < \omega + k$, which is absurd since ordinals are their own ranks.
\end{boxnexample}

Let's make some definitions about cardinals.

\begin{boxdefinition}[(Strong) Limit Cardinals]
    $\mu$ is a
    \begin{itemize}
        \item \textbf{limit cardinal} iff for all cardinals $\kappa < \mu$, $\kappa^+ < \mu$.
        \item \textbf{strong limit cardinal} iff for all cardinals $\kappa < \mu$, $2^{\kappa} < \mu$.
    \end{itemize}
\end{boxdefinition}

\begin{boxdefinition}[Cofinality]
    % [CORRECTED] was: ``of a set $A$'' --- ``unbounded in $A$'' needs an order on $A$, and the notes only take cofinalities of ordinals
    The \textbf{cofinality} of an ordinal $A$, denoted $\cf(A)$, is the least $\abs{B}$ such that there is a function $f : B \to A$ whose image $f[B]$ is unbounded in $A$.
\end{boxdefinition}

\begin{boxdefinition}[Regularity/Singularity]
    A cardinal $\mu$ is \textbf{regular} if $\cf(\mu) = \mu$. If $\mu$ is not regular, we say $\mu$ is \textbf{singular}.
\end{boxdefinition}

\begin{boxdefinition}[Weak/Strong Inaccessibility]
    We say a cardinal $\mu$ is
    \begin{itemize}
        \item \textbf{weakly inaccessible} iff $\mu$ is an uncountable, regular limit cardinal.
        \item \textbf{strongly inaccessible} iff $\mu$ is an uncountable, regular, strong limit cardinal.
    \end{itemize}
\end{boxdefinition}

\begin{boxproposition}
    If $\mu$ is strongly inaccessible, then $\beth_{\mu} = \mu$.
\end{boxproposition}
\begin{proof}
    By induction, $\beth_{\alpha} \geq \alpha$ for every ordinal $\alpha$, so $\beth_{\mu} \geq \mu$.

    It's enough to show that $\forall \alpha < \mu$, $\beth_{\alpha} < \mu$, because the fact that $\mu$ is a limit ordinal means that, by definition, $\beth_{\mu} = \sup_{\alpha < \mu} \beth_{\alpha}$.

    Now perform induction on $\alpha < \mu$.
    
    Base case: $\beth_0 < \mu$ because $\beth_0 = \omega$ and $\mu$ is uncountable.

    Successor case: if $\alpha + 1 < \mu$ and $\beth_{\alpha} < \mu$, then $\beth_{\alpha + 1} = 2^{\beth_{\alpha}} < \mu$, because $\beth_{\alpha}$ is a cardinal below $\mu$ and $\mu$ is a strong limit.

    Limit case: if $\alpha < \mu$ is a limit and $\beth_{\gamma} < \mu$ for all $\gamma < \alpha$, then $\beth_{\alpha} = \sup_{\gamma < \alpha} \beth_{\gamma}$ is a supremum of $\abs{\alpha} < \mu$ many ordinals below $\mu$, so it is below $\mu$, as $\mu$ is regular.
\end{proof}

% 23 September

We now continue going through the axioms.

Let's look at Foundation.

We want to see if $\F^M$, that is,
\begin{align*}
    \brac{\forall A \parenth{A \ne \emptyset \to \exists x \in A, \, \forall y \in A, \, y \notin x}}^M
\end{align*}
Consider the statement
\begin{align*}
    \forall A \in M \parenth{A \ne \emptyset \to \exists x \in A, \, \forall y \in A, \, y \notin x}
\end{align*}
This is true if $M \ssq \WF$.

We know the following:
\begin{itemize}
    \item if $A \in \WF$, then $A \subset \WF$
    \item if $x \in A$ with $\rank{x} = \min\setst{\rank{y}}{y \in A}$ then $x$ is $\in$-minimal in $A$.
    \item if $y \in x$ then $\rank{y} < \rank{x}$
\end{itemize}
In particular, $\F^{\WF}$ and $\F^{V_{\alpha}}$.

Finally, for infinity, we can show that if $\omega \in M$, then infinity relativises to $M$; in particular, infinity relativises to $\WF$ and to $V_{\alpha}$ for $\alpha > \omega$ a limit.

We're now done showing that every axiom of $\ZF$ without infinity relativises to $V_{\omega}$, and furthermore, that every axiom of $\ZF$ relativises to $\WF$. As we saw earlier, this gives us the following.

\begin{boxcorollary}
If $\ZFwoF$ is consistent, then so is $\ZF$.
\end{boxcorollary}

Now, let's study choice. We know that under $\ZFwoF$, this is equivalent to the well-ordering principle.

We want to show, effectively, that
\begin{align*}
    \brac{\forall A \exists R \parenth{(A, R) \text{ is a w.o.}}}^M
\end{align*}
% [CORRECTED] was: ``This is equivalent to $\forall A \in M, \, \exists R \in M \, \brac{\parenth{(A, R) \text{ is a w.o.}}}^M$ because \ldots'' --- that equivalence holds by the definition of relativization, so the reason given did no work; the absoluteness is what transfers ``w.o.'' from $V$ down to $M$, which is how the argument below uses it
This follows from
\begin{align*}
    \forall A \in M, \, \exists R \in M \, \parenth{(A, R) \text{ is a w.o.}}
\end{align*}
because the statement ``$(A, R)$ is a w.o.'' is $\Pi_1$, and thus, downwards absolute.

Indeed, the reason this statement is $\Pi_1$ is that $(A, R)$ is a w.o. iff $(A, R)$ is a l.o. (which is a $\Delta_0$ statement) and $(A, R)$ is well-founded (which is $\Pi_1$ because it says $\forall S \brac{\parenth{\emptyset \ne S \ssq A} \to \exists x \in S \, \forall y \in S \, \neg\parenth{y \, R \, x}}$, and everything after $\forall S$ is $\Delta_0$: the quantifiers over $x$ and $y$ are bounded by $S$, and $\neg\parenth{y \, R \, x}$ says that the ordered pair $\parenth{y, x}$ is not a member of $R$, which needs only quantifiers bounded by $R$). % [FILLED] the reason well-foundedness is Pi_1 is supplied here

In the background, now, assume $\ZFCwoF$. Let $A \in \WF$ and suppose $A \ssq V_{\alpha}$ for some $\alpha$. In $V$, we have some $R \ssq A \times A$ such that $R$ is a w.o. Then $R \ssq V_{\alpha} \times V_{\alpha}$, so $R \in V_{\alpha + 3} \subset \WF$. By $\Pi_1$ downwards absoluteness, we conclude that $\brac{(A, R) \text{ is a w.o.}}^{\WF}$.

\begin{boxcorollary}
    $(\mathbf{AC})^{\WF}$, $(\mathbf{AC})^{\V_{\alpha}}$ for all limit $\alpha$.
\end{boxcorollary}

\begin{boxcorollary}
    If $\ZFCwoF$ is consistent, then so is $\ZFC$.
\end{boxcorollary}

\begin{boxcorollary}
    If $M$ is strongly inaccessible, then $\sigma^{V_M}$ for all axioms $\sigma$ of $\ZFC$.
\end{boxcorollary}

\begin{boxlemma}
    % [CORRECTED] was: ``without power sets, infinity and foundation'' --- the proof needs Foundation to make ``$x$ is an ordinal'' a $\Delta_0$ condition, and without Foundation the statement fails (a transitive, $\in$-linear, ill-founded set can witness the $\Sigma_1$ form)
    If we denote by $T$ the theory consisting of $\ZF$ without power sets and infinity, the statement ``$(A, R)$ is a w.o.'' is also $\Sigma_1^T$.
\end{boxlemma}
\begin{proof}
    $(A, R)$ is a w.o. iff there is some $F : A \to \OR$ such that $\forall x, y \in A \parenth{x \, R \, y \to F(x) < F(y)}$ % [FILLED] the rest of the proof is supplied here; the lecture stated the characterization and stopped
    and $(A, R)$ is a l.o. From left to right, take $F$ to be the collapsing map of \Cref{Ch1:Thm:Mostowski_WO}, which is order-preserving with ordinal values; its construction is a recursion along a well-ordering of a set, and uses Replacement but neither Power Set nor Infinity. From right to left, given a nonempty $S \ssq A$, the set $F[S]$ is a nonempty set of ordinals, so it has a least element $F(x)$ with $x \in S$, and no $y \in S$ has $y \, R \, x$, as that would give $F(y) < F(x)$.

    So, in $T$, ``$(A, R)$ is a w.o.'' is equivalent to
    \begin{align*}
        \text{$(A, R)$ is a l.o.} \land \exists F \brac{F \text{ is a function} \land \dom{F} = A \land \forall x \in A \, \parenth{F(x) \in \OR} \land \forall x, y \in A \, \parenth{x \, R \, y \to F(x) \in F(y)}}
    \end{align*}
    Everything after $\exists F$ is $\Delta_0$ provided ``$F(x) \in \OR$'' is, and this is where Foundation is used: in its presence, an ordinal is just a transitive set that is linearly ordered by $\in$, which is a $\Delta_0$ condition.
\end{proof}

% 25 September

For today, let $M$ be a transitive class model of $\ZFwoP$.

\begin{boxlemma}
    $[M]^{< \omega} \ssq M$.
\end{boxlemma}
\begin{proof}
    Given $x_1. \ldots, x_n \in M$, $\set{x_1, \ldots, x_n} = \bigcup_{j=1}^{n} \set{x_j}$.
\end{proof}

The formula ``$A$ is finite'' is absolute for $M$. In fact, we can do better.

\begin{boxlemma}
    For each $n < \omega$, the formula ``$\abs{A} = n$'' is absolute for $M$.
\end{boxlemma}
\begin{proof}
    \hfill 
    \begin{description}
        \item[\underline{Upwards absoluteness.}] Say $f \in M$ and
        \begin{align*}
            \parenth{f : n \iso A}^M
        \end{align*}
        where $f : n \iso A$ means $f$ is injective and surjective.

        We know that $f : n \iso A$ because we already know it is absolute.

        \item[\underline{Downwards absoluteness.}] We have that $f : n \iso A$ in $V$. But $f \ssq_{\text{finite}} n \times A \ssq M$. And indeed we've checked already that Cartesian products are absolute.

        Then, by the previous lemma, $f \in M$.

        Now again, by $\Delta_0$-absoluteness,
        \begin{align*}
            \parenth{f : n \iso A}^M
        \end{align*}
    \end{description}
\end{proof}

\begin{boxlemma}
    ``$A$ infinite'' is absolute for $M$.
\end{boxlemma}
\begin{proof}
    Clear from previous.
\end{proof}

\begin{boxlemma}
    If $A \in M$ and $A$ is finite (in both $M$ and $V$), then $\Powset(A) \in M$. Therefore, $A \mapsto \Powset(A)$, restricted to finite sets $A$, is an absolute class function for $M$. % [CORRECTED] ``restricted to finite sets $A$'' inserted --- for infinite $A$ it fails, eg, $\Powset(\omega) \notin \HC$

    For the remainder of this argument, let's work in $M$.
    
    We argue by recursion. We have that $A$ is finite, so there is some $n < \omega$ and $f : n \iso A$. Denote by $2^n$ the set of all $n$-tuples of $0$s and $1$s. For each $x \in 2^n$, write
    \begin{align*}
        B_x = \setst{f(i)}{x_i = 1} \ssq A
    \end{align*}
    In particular, we have a mapping $x \mapsto B_x$. Finally, we can conclude that
    \begin{align*}
        \Powset(A) = \setst{B_x}{x \in 2^n}
    \end{align*}
\end{boxlemma}

\subsection{Absoluteness of Recursive Definitions}

Often, we have $(A, R)$ well-founded and set-like, and some $F : A \times V \to V$. We then define $G : A \to V$ by recursion so that
\begin{align*}
    \forall y \in A, \, G(y) = F\of{y, G \restriction \setst{x}{x \, R \, y}}
\end{align*}
Now, we'd like to know that if the ``data'' above is absolute, then so is $G$. How do we show this?

Ok, let's list our assumptions (recalling that a class function is absolute iff the defining formula is absolute, and also $M$ is closed under the values):
\begin{itemize}
    \item $A$ and $R$ are absolute
    \item $F$ is an absolute class function
    \item $y \mapsto \setst{x}{x \, R \, y}$ is an absolute class function
\end{itemize}
We can conclude from this the following;
\begin{itemize}
    \item $\brac{\text{$(A, R)$ is set-like and well-founded}}^M$
    \item $G$ is absolute for $M$
    \item In particular, for every $y \in A^M$,
    \begin{align*}
        G^M(y) = F^M\of{y, G^M \restriction \setst{x}{x \, R^M y}}
    \end{align*}
    and
    \begin{align*}
        G^M = G \restriction A^M = G \restriction \parenth{A \cap M}
    \end{align*}
\end{itemize}

\begin{boxlemma}
    $\brac{\text{$(A, R)$ is set-like and well-founded}}^M$
\end{boxlemma}
\begin{proof}[Proof sketch]
    Note that $(A, R)$ is well-founded iff
    \begin{align*}
        \forall S, \, \brac{\emptyset \ne S \ssq A \to \exists x \in S, \, \forall y \in S,\, \neg \parenth{y \, R \, x}}
    \end{align*}
    We can essentially rewrite the expression inside brackets using quantifiers that are bounded over $(A, R)$ (subject to the universal $\forall S$). The bracketed expression is absolute between $M$ and $V$, and the universal out front makes it downwards absolute from $V$ to $M$.

    Note that set-like is downwards absolute from $V$ to $M$ because we assumed that $y \mapsto \setst{x}{x \, R \, y}$ and $(A, R)$ are all absolute.
\end{proof}

\begin{boxlemma}
    $G$ is absolute for $M$. In particular, for every $y \in A^M$,
    \begin{align*}
        G^M(y) = F^M\of{y, G^M \restriction \setst{x}{x \, R^M y}}
    \end{align*}
    and
    \begin{align*}
        G^M = G \restriction A^M = G \restriction \parenth{A \cap M}
    \end{align*}
\end{boxlemma}
\begin{proof}[Proof sketch]
    Just show that for all $y$ in $A^M = A \cap M$, we have $G^M(y) = G(y)$ by induction on $y$.
\end{proof}

There are many useful applications of absoluteness.

\begin{boxexample}[Defining Ordinal Arithmetic]
    There are many ways we can use absoluteness to define ordinal arithmetic.
    \begin{itemize}
        \item \underline{Defining it recursively.} For example, we can define $\alpha + 0 = \alpha$, $\alpha + \parenth{\beta + 1} = \parenth{\alpha + \beta} + 1$, and if $\gamma$ is a limit ordinal, we can define $\alpha + \gamma := \sup_{\beta < \gamma}\of{\alpha + \beta}$.

        \item \underline{Defining it using Mostowski Collapse of well-orderings.} For example, we can define
        \begin{align*}
            \alpha + \beta :=
            \type\of{\parenth{\parenth{\set{0} \times \alpha} \cup \parenth{\set{1} \times \beta}, <_{\lex}}}
        \end{align*}
        Note that $\type$ is the range of the Mostowski Collapse. Indeed,
        \begin{align*}
            G(y) = \setst{G(x)}{x <_{\lex} y}
        \end{align*}
        is absolute.
    \end{itemize}
\end{boxexample}

% One final remark for today. We know that $M$ is a transitive class. If $M$ is a proper class, then $\OR^M = \OR$. If $M$ is a set, then $\OR^M \in \OR$ (ie, if $M$ is a set, then as a set of ordinals that is transitive, it must itself be a 

% 28 Sept

If $M$ is a transitive model of $\ZFwoP$, we claim that the statement ``$\alpha$ is an ordinal'' is absolute for $M$. % [CORRECTED] was: ``$\alpha$ is absolute for $M$'' --- absoluteness is a property of formulae, and what the next display unpacks is the formula ``$\alpha$ is an ordinal'\,'

But what does it even mean? It means
\begin{align*}
    \OR^M = \OR \cap M
\end{align*}
As $M$ is transitive, $\OR^M$ is transitive, so either $\OR^M \in \OR$ or $\OR^M = \OR$. It turns out:
\begin{enumerate}
    \item $\OR^M \in \OR$ iff $M$ is a set
    \item $\OR^M = \OR$ iff $M$ is a proper class
\end{enumerate}
To see this, note that foundation is equivalent to $(V, \in)$ being well-founded. It's clear that $(V, \in)$ is set-like: if $y$ is a set, then everything that is $\in$-below $y$ is just $y$, which is a set. So by the Theorem Scheme on Recursion (which is true in $\ZFwoP$), we have a class function $G : V \to \OR$ such that $\forall y$, $G(y) = \sup_{x \in y}\of{G(x) + 1}$.

Now if we also give ourselves the power set axiom, then we can define the $V$-hierarchy, and conclude that $G(y) = \rank{y}$. But without power sets, \textit{we cannot even define the $V$-hierarchy}, so there is no meaningful notion of rank!

By the theorem scheme on absoluteness of recursively defined class functions, $G$ is an absolute class function for $M$. In other words, $G^M = G \restriction M$. Also, $G^M : M \to \OR^M$ is surjective. % [FILLED] the surjectivity and the two cases of the dichotomy are argued here; the lecture asserted them
Indeed, each value of $G^M$ is an ordinal in $M$, and $G(\alpha) = \alpha$ for every ordinal $\alpha$, by induction: $G(\alpha) = \sup_{\beta < \alpha}\of{G(\beta) + 1} = \sup_{\beta < \alpha}\of{\beta + 1} = \alpha$. So every $\alpha \in \OR^M = \OR \cap M$ is $G^M(\alpha)$.

This gives both halves of the dichotomy.
\begin{itemize}
    \item If $M$ is a set, then $\OR \cap M = G[M]$ is a set by Replacement, and a transitive set of ordinals is an ordinal, so $\OR^M \in \OR$.
    \item If $M$ is a proper class, suppose $\OR \cap M$ were a set, hence an ordinal $\alpha$. In $V$ we do have Power Set, so $G(y) = \rank{y}$ for every $y$. So, for every $y \in M$, $\rank{y} = G^M(y) \in \OR \cap M = \alpha$, ie, $y \in V_{\alpha}$. Then $M \ssq V_{\alpha}$, and $M$ would be a set by Comprehension. So $\OR^M$ is a transitive proper class of ordinals, which can only be $\OR$.
\end{itemize}

Now let's study the uncountable ordinals in $M$.

Denote by $\omega_1^M$ what ``$M$ thinks'' is the least uncountable ordinal. If $\parenth{\omega_1 \text{ exists}}^M$, ie, if the notion of the existence of the least uncountable ordinal is even definable in $M$, then $\omega_1^M \leq \omega_1$ (where $\le$ is just the usual $\le$ on ordinals---recall that $\OR$ is absolute, which means precisely that whatever a class \textit{thinks} is an ordinal is, indeed, an ordinal in $V$).

\begin{boxdefinition}
    For each cardinal $\kappa \geq \omega$, define the class
    \begin{align*}
        H_{\kappa} = \setst{z}{\abs{\trcl{\set{z}}} < \kappa}
    \end{align*}
\end{boxdefinition}

Note that
\begin{align*}
    \trcl{\set{z}} = \set{z} \cup \trcl{z}
\end{align*}
so
\begin{align*}
    \abs{\trcl{\set{z}}} = 1 \+ \abs{\trcl{z}}
\end{align*}

We can also unwind $\trcl{z}$ one level further, through the members of $z$.

% [CORRECTED] was: $\bigcup_{x \in z} \trcl{\set{z}}$ --- true, but only because the index $x$ never appears, so it says nothing; the uses below need $\trcl{\set{x}}$
\begin{boxproposition}
    For all sets $z$,
    \begin{align*}
        \trcl{\set{z}} = \set{z} \cup \bigcup_{x \in z} \trcl{\set{x}}
    \end{align*}
\end{boxproposition}

Nothing is assumed about $z$ here: it need not be transitive, let alone an ordinal.

We will work in $\ZFC$ because we need choice to be able to work comfortably with cardinalities.

\begin{boxlemma}
    Each $H_{\kappa}$ is a transitive class and, if $\kappa < \lambda$, then $H_{\kappa} \subsetneq H_{\lambda}$.
\end{boxlemma}
\begin{proof}
    \hfill
    \begin{description}
        \item[\underline{Transitivity.}] Let $z \in H_{\kappa}$ and $y \in z$. Then $y \in \trcl{\set{z}}$, so $\trcl{\set{z}}$ is a transitive set with $y$ as a member, and so it contains the smallest such set: $\trcl{\set{y}} \ssq \trcl{\set{z}}$. Hence $\abs{\trcl{\set{y}}} \leq \abs{\trcl{\set{z}}} < \kappa$, and $y \in H_{\kappa}$.

        \item[\underline{Monotonicity.}] If $\kappa < \lambda$, then $\abs{\trcl{\set{z}}} < \kappa$ implies $\abs{\trcl{\set{z}}} < \lambda$, so $H_{\kappa} \ssq H_{\lambda}$. The inclusion is strict because $\kappa \in H_{\lambda} \setminus H_{\kappa}$: $\trcl{\set{\kappa}} = \kappa + 1$ (it is transitive and has $\kappa$ as a member, and any transitive set with $\kappa$ as a member contains $\kappa$ and all its members), and $\abs{\kappa + 1} = \kappa$, as $\kappa$ is infinite, which is below $\lambda$ but not below $\kappa$.
    \end{description}
\end{proof}

\begin{boxlemma}
    For all $\kappa$, $H_{\kappa} \subseteq V_{\kappa}$.
\end{boxlemma}
\begin{proof}
    Let $z \in H_{\kappa}$. We will show that $z \in V_{\kappa}$. Let
    \begin{align*}
        \sigma = \setst{\rank{y}}{y \in \trcl{\set{z}}}
    \end{align*}
    Then $\sigma$ is a set of ordinals and $\rank{z} = \max(\sigma)$.

    We claim that $\sigma$ is a transitive set. Indeed, if not, then we would have some $\alpha < \beta$ such that $\beta = \min\of{\sigma \setminus \alpha}$. Pick $y \in \trcl{\set{z}}$ such that $\rank{y} = \beta$. Then, $\forall x \in y$, since $x \in \trcl{\set{y}}$ and $\rank{x} < \rank{y}$, we conclude that $\rank{x} < \alpha$. Then, by the rank formula,
    \begin{align*}
        \beta = \rank{y} = \sup_{x \in y}\of{\rank{x} + 1} \leq \alpha
    \end{align*}
    This is a contradiction! So $\sigma$ is transitive as desired.

    This then tells us that $\sigma = \rank{z} + 1$: a transitive set of ordinals is an ordinal, and an ordinal whose largest element is $\rank{z}$ is $\rank{z} + 1$.
    Then, we get
    \begin{align*}
        \abs{\sigma} \leq \sigma < \abs{\sigma}^+ \leq \abs{\trcl{\set{z}}}^+ \leq \kappa
    \end{align*}
    so $\rank{z} < \kappa$.
\end{proof}

\begin{boxcorollary}
    Each $H_{\kappa}$ is a set.
\end{boxcorollary}
\begin{proof}
    Just apply comprehension.
\end{proof}

\begin{boxdefinition}
    Define $\HF := H_{\omega}$ and $\HC := H_{\omega_1}$
\end{boxdefinition}

We will show that $\HC$ is a model of $\ZFwoP$. In fact, so is $H_{\kappa^+}$ for each cardinal $\kappa$.

\begin{boxlemma}
    $\HF = V_{\omega}$
\end{boxlemma}
\begin{proof}
    We just saw that $\HF \ssq V_{\omega}$. Given some $z \in V_{\omega}$, pick $n < \omega$ such that $z \in V_n$. Then
    \begin{align*}
        \abs{\trcl{\set{z}}} \leq \abs{V_n} < \omega
    \end{align*}
    so $z \in \HF$.
\end{proof}

\begin{boxlemma}
    $\OR \cap H_{\kappa} = \kappa$
\end{boxlemma}
\begin{proof}
    If $\alpha \in \OR$, then $\trcl{\set{\alpha}} = \alpha + 1$, because $\alpha + 1 = \alpha \cup \set{\alpha}$ is transitive and has $\alpha$ as a member, and any transitive set with $\alpha$ as a member contains $\alpha$ and all its members. So $\alpha \in H_{\kappa}$ iff $\abs{\alpha + 1} < \kappa$.

    If $\alpha < \kappa$, then $\alpha + 1 < \kappa$ too, because $\kappa$ is an infinite cardinal and so a limit ordinal; so $\abs{\alpha + 1} \leq \alpha + 1 < \kappa$. If $\alpha \geq \kappa$, then $\abs{\alpha + 1} \geq \abs{\kappa} = \kappa$. So the ordinals in $H_{\kappa}$ are exactly those below $\kappa$.
\end{proof}

\begin{boxlemma}
    If $H_{\kappa} = V_{\alpha}$, then $\kappa = \alpha$. 
\end{boxlemma}

% 30 September

Recall that $\setst{\kappa}{\kappa = \beth_{\kappa}}$ is a club class of cardinals. It is closed because the $\beth$ operation is ``continuous'': if $\lambda$ is a limit of $\beth$-fixed points, then, as $\alpha \mapsto \beth_{\alpha}$ is increasing, $\beth_{\lambda} = \sup_{\alpha < \lambda} \beth_{\alpha}$ is also the supremum of $\beth_{\kappa} = \kappa$ over the fixed points $\kappa < \lambda$, which is $\lambda$. So a sup of $\beth$-fixed points is a $\beth$-fixed point. It is unbounded because, given any $\alpha$, the sequence $\kappa_0 = \alpha$, $\kappa_{n + 1} = \beth_{\kappa_n}$ is non-decreasing (as $\beth_{\gamma} \geq \gamma$) and its supremum $\lambda \geq \alpha$ is a fixed point: either the sequence is eventually constant, at a fixed point, or $\lambda$ is a limit and $\beth_{\lambda} = \sup_{n} \beth_{\kappa_n} = \sup_n \kappa_{n + 1} = \lambda$.

We also saw that if $\mu$ is a strongly inaccessible cardinal, then $\mu = \beth_{\mu}$.

Here is a useful lemma.

\begin{boxlemma}
    If $\kappa$ is regular, then
    \begin{align*}
        H_{\kappa} = \setst{z}{\forall x \in \trcl{\set{z}}, \, \abs{x} < \kappa}
    \end{align*}
\end{boxlemma}
\begin{proof}
    Call the RHS $I$. We show $H_{\kappa} = I$.

    \begin{description}
        \item[\underline{$H_{\kappa} \ssq I$.}] Let $z \in H_{\kappa}$ and $x \in \trcl{\set{z}}$. Then $x \ssq \trcl{\set{z}}$ and $\abs{x} \le \abs{\trcl{\set{z}}} < \kappa$.

        \item[\underline{$H_{\kappa} \supseteq I$.}] We apply $\in$-induction. Ie, we assume that for $x \in z$, since $x \in I$ (by transitivity of $I$), $x \in H_{\kappa}$, and therefore, $\abs{\trcl{\set{x}}} < \kappa$.

        Performing the induction, we calculate
        \begin{align*}
            \abs{\trcl{\set{z}}} \le 1 \+ \bigoplus_{x \in z} \underbrace{\abs{\trcl{\set{x}}}}_{< \kappa}
        \end{align*}
        As $z \in I$, $\abs{z} < \kappa$. So we are summing $< \kappa$ many cardinals that are all $< \kappa$, resulting in something $< \kappa$.
    \end{description}
\end{proof}

So now let's ask ourselves what Axioms of ZFC hold in these $H_{\kappa}$.

\begin{boxlemma}
    The following ZFC axioms hold in all $H_{\kappa}$:
    \begin{enumerate}
        \item Extensionality
        \item Comprehension
        \item Pairing
        \item Union
        \item Choice
        \item Foundation \textit{(if we assume foundation in $V$)}
    \end{enumerate}
\end{boxlemma}
\begin{proof} \hfill
    \begin{enumerate}
        \item $H_{\kappa}$ is transitive.
        \item If $A \in H_{\kappa}$ and $B \ssq A$, then $B \in H_{\kappa}$.
        % [CORRECTED] items 3 and 4 were both: ``Essentially just because $H_{\kappa}$ is a set.'' --- that a class is a set gives no member of it containing $x$ and $y$ (eg, $2 = \set{0, 1}$ is a transitive set in which Pairing fails); what is needed is closure of $H_{\kappa}$ under pairs and unions
        \item If $x, y \in H_{\kappa}$, then $\set{x, y} \in H_{\kappa}$, since $\trcl{\set{\set{x, y}}} = \set{\set{x, y}} \cup \trcl{\set{x}} \cup \trcl{\set{y}}$ has cardinality below $\kappa$, as $\kappa$ is infinite.
        \item If $A \in H_{\kappa}$, then $\bigcup A \in H_{\kappa}$, since $\bigcup A \ssq \trcl{A}$, and so $\trcl{\set{\bigcup A}} \ssq \set{\bigcup A} \cup \trcl{A}$.
        \item Let $A \in H_{\kappa}$. Pick a well-ordering $R$ of $A$ (assuming AC in $V$). Then, as $R \ssq A \times A$, $R \in H_{\kappa}$. Use downwards $\Pi_1$-absoluteness to see that
        \begin{align*}
            \brac{(A, R) \text{ is a well-ordering}}^{H_{\kappa}}
        \end{align*}
    \end{enumerate}
\end{proof}

\begin{boxlemma}
    Infinity holds in $H_{\kappa}$ iff $\kappa > \omega$.
\end{boxlemma}
\begin{proof}
    Follows from the fact that $\kappa > \omega \lr \omega \in H_{\kappa}$.
\end{proof}

\begin{boxlemma} \label{Ch3:Lemma:H_Kappa_Replacement}
    If $\kappa = \cf(\kappa)$ (ie, $\kappa$ is regular), then each axiom $\sigma$ of the Replacement scheme holds relativised to $H_{\kappa}$, ie, $\sigma^{H_{\kappa}}$.
\end{boxlemma}
\begin{proof}
    If $A \in H_{\kappa}$ and $f : A \to H_{\kappa}$, then, letting $B = \range\of{f}$, we have
    \begin{align*}
        \abs{\trcl{\set{B}}}
        &= \abs{\set{B} \cup \bigcup_{x \in A} \trcl{\set{f(x)}}} \\
        &\leq 1 \+ \underbrace{\sum_{x \in A}}_{\text{fewer than } \kappa \text{ terms}} \underbrace{\abs{\trcl{\set{f(x)}}}}_{\text{each } < \kappa} \\
        &< \kappa
    \end{align*}
    There are strictly fewer than $\kappa$ terms because $A \ssq \trcl{\set{A}}$, and a sum of fewer than $\kappa$ cardinals, each below $\kappa$, is below $\kappa$ because $\kappa$ is regular. So $B \in H_{\kappa}$.

    This is exactly what Replacement in $H_{\kappa}$ needs. If $A, \bar{z} \in H_{\kappa}$ and each $x \in A$ has a unique $y \in H_{\kappa}$ with $\vp^{H_{\kappa}}\of{x, y, \bar{z}}$, let $f(x)$ be that $y$; $f$ is a set by Replacement in $V$, as $A \ssq H_{\kappa}$, and $B = \range\of{f} \in H_{\kappa}$ is the witness the axiom asks for.
\end{proof}

\begin{boxlemma}
    If $\lambda = \kappa^+$, then
    \begin{align*}
        \parenth{\neg \text{(Power Set)}}^{H_{\lambda}}
    \end{align*}
\end{boxlemma}
\begin{proof}
    $\Powset(\kappa) \ssq H_{\lambda}$ but $\Powset(\kappa) \notin H_{\lambda}$.
\end{proof}

\begin{boxlemma} \label{Ch3:Lemma:H_Kappa_Power_Set}
    TFAE:
    \begin{enumerate}[label = (\arabic*)]
        \item $(\text{Power Set})^{H_{\lambda}}$
        \item $\lambda$ is a strong limit
    \end{enumerate}
\end{boxlemma}
\begin{proof} \hfill
    \begin{description}
        \item[\underline{$(2) \to (1)$.}] Assume $(2)$. Let $A \in H_{\lambda}$. Then
        \begin{align*}
            \abs{\trcl{\set{\Powset(A)}}}
            = \abs{\set{\Powset(A)} \cup \Powset(A) \cup \trcl{A}}
            \leq 1 \+ 2^{\abs{A}} \+ \abs{\trcl{A}}
            < \lambda
        \end{align*}

        \item[\underline{$\neg(2) \to \neg(1)$.}] We show, effectively, that if $\kappa < \lambda$ is a cardinal with $2^{\kappa} \geq \lambda$, then $\kappa \in H_{\lambda}$ and $\Powset(\kappa) \ssq H_{\lambda}$ but $\Powset(\kappa) \notin H_{\lambda}$, so Power Set fails in $H_{\lambda}$.
    \end{description}
\end{proof}

\begin{boxcorollary} \hfill
    \begin{enumerate}
        \item $H_{\omega} = V_{\omega}$ models $\ZFCwoI$
        \item $H_{\omega_1} = H_{\omega_1}$ and all other $H_{\omega_2}, H_{\omega_3}, \ldots$ model $\ZFCwoP$
    \end{enumerate}
\end{boxcorollary}

The first cardinal not covered by the corollary is $\aleph_{\omega}$, which is singular, and it shows that the regularity hypothesis in \Cref{Ch3:Lemma:H_Kappa_Replacement} cannot be dropped.

\begin{boxexample}[$H_{\aleph_{\omega}}$]
    Everything but Replacement and Power Set holds in $H_{\aleph_{\omega}}$: Extensionality, Comprehension, Pairing, Union, Choice and Foundation hold in every $H_{\kappa}$, and Infinity holds because $\aleph_{\omega} > \omega$.

    Replacement fails. There's a formula $\vp\of{u, v}$ so that, for all $n \in \omega$ and $\kappa \in H_{\aleph_{\omega}}$,
    \begin{align*}
        \vp\of{n, \kappa}^{H_{\aleph_{\omega}}} \lr \kappa = \aleph_n
    \end{align*}
    Namely, $\vp\of{n, \kappa}$ says that there is a function $g$ with domain $n + 1$ such that $g(0) = \omega$, each $g(i + 1)$ is the least cardinal above $g(i)$, and $g(n) = \kappa$. This relativizes correctly because the witness $\grpres{\aleph_i}{i \leq n}$ lies in $H_{\aleph_{\omega}}$, and because an ordinal below $\aleph_{\omega}$ is a cardinal in $H_{\aleph_{\omega}}$ iff it is one in $V$, as every bijection between two ordinals below $\aleph_{\omega}$ lies in $H_{\aleph_{\omega}}$. So $\vp$ defines $f : n \mapsto \aleph_n$ over $H_{\aleph_{\omega}}$, on the domain $\omega \in H_{\aleph_{\omega}}$, but $f \notin H_{\aleph_{\omega}}$.
    \begin{figure}[H]
        \centering
        \begin{tikzpicture}[dot/.style = {circle, fill, inner sep = 1.5pt}, every node/.style = {font = \small}]
            \drawcone{2}{4.5}{$H_{\aleph_{\omega}}$}
            \node[dot, label = {left:$\omega$}] (w) at (0.1, 0.7) {};
            \node[dot, label = {left:$\aleph_n$}] (a) at (-0.5, 3.2) {};
            \draw[->, thick] (w) to[bend right = 60] node[right] {$f : n \mapsto \aleph_n$} (a);
        \end{tikzpicture}
    \end{figure}
    Indeed, any $B$ with every $\aleph_n \in B$ has $\aleph_{\omega} = \bigcup_{n < \omega} \aleph_n \ssq \trcl{\set{B}}$, so $\abs{\trcl{\set{B}}} \geq \aleph_{\omega}$ and $B \notin H_{\aleph_{\omega}}$. So no set in $H_{\aleph_{\omega}}$ collects the values of $f$.

    Whether Power Set holds in $H_{\aleph_{\omega}}$ is, by \Cref{Ch3:Lemma:H_Kappa_Power_Set}, the question of whether $\aleph_{\omega}$ is a strong limit, which $\ZFC$ does not settle.
\end{boxexample}

\begin{boxlemma}
    Let $\kappa$ be an uncountable cardinal. Then $H_{\kappa} \prec_{\Sigma_1} V$.
\end{boxlemma}
\begin{remark}
    This means that for each $\Sigma_1$ formula $\vp\of{\bar{u}}$ and each matching $\bar{x} \in H_{\kappa}$, $\vp\of{\bar{x}}^{H_{\kappa}} \lr \vp\of{\bar{x}}^V$. In other words, it means $\Sigma_1$-formulae are absolute for $H_{\kappa}$. We also say, ``$H_{\kappa}$ is $\Sigma_1$-correct in $V$.''
\end{remark}
\begin{proof}
    We already know upwards absoluteness, so the only remaining question is downwards absoluteness.

    % [FILLED] forward reference supplied here; the lecture deferred this proof until satisfaction had been formalized
    This needs us to talk about satisfaction in sets, which we have not yet formalized; we do so in the next section, and the proof is given there (\Cref{Ch3:Thm:H_Kappa_Sigma1}).
\end{proof}

\section{Relativisaiton vs Satisfaction}

Can we formalise into Set Theory the notion of ``$M \models \vp\of{\bar{x}}$''?

Pick ``symbols'' and formalise the following using $\ZFwoP$:
\begin{itemize}
    \item syntax of first-order logic
    \item definition of satisfaction
    \item proof theory
    \item completeness theorem
\end{itemize}

The way we define what ``formulae'' are internally is using Gödel numbers. Ie, for each formlua $\vp$ (\textit{external}), there exists some set $\ulcorner \vp \urcorner$ that is analogous to it. We still write $M \models \vp\of{\bar{x}}$.

It is possible to show that our analogies translate both ways, ie, $\vp\of{\bar{x}}^M$ iff ``$M \models \vp\of{\bar{x}}$'', simply by structural induction on $\vp$.

% Transcribed from the board photographs supplied with the post-lecture pass of 8 Oct 2026; these boards were not typed up during the lecture.

Note, about $M \models \vp\of{\bar{x}}$:
\begin{enumerate}
    \item $\vp$ is not a set. There's a set $\ulcorner \vp \urcorner$ that is analogous to $\vp$. Nevertheless, we just write $M \models \vp\of{\bar{x}}$.
    \item This requires $M \ne \emptyset$ and $\bar{x} \in M^{n + 1}$, where $v_0, \ldots, v_n$ are the free variables of $\vp$. Contrast with $\vp^M\of{\bar{x}}$.
    \item $\vp^M\of{\bar{x}}$ was defined for each $\vp$ and each class $M$. $M \models \vp\of{\bar{x}}$ is defined only for sets $M$:
    \begin{align*}
        \models \; \ssq \; \setst{\parenth{M, c, \bar{x}}}{M \ne \emptyset, \, c \text{ a Gödel number}, \, \bar{x} \in M^{<\omega}}
    \end{align*}
\end{enumerate}

\begin{boxlemma}[Fact (Scheme)]
    Consider any formula $\vp\of{v_0, \ldots, v_n}$. For each $M \ne \emptyset$ and $\bar{x} \in M^{n + 1}$,
    \begin{align*}
        \vp\of{\bar{x}}^M \text{ iff } ``M \models \vp\of{\bar{x}}"
    \end{align*}
\end{boxlemma}

For now, work in $\ZF$. Given a formula $\vp(u)$ and a set $x$, say $\vp(u)$ \textbf{defines} $x$ iff
\begin{align*}
    \forall u \brac{u = x \lr \vp(u)} \tag{$\star$}
\end{align*}
Observe: there is no formula $\delta\of{r, x}$ so that, for every formula $\vp(u)$,
\begin{align*}
    \forall x \brac{\delta\of{\ulcorner \vp(u) \urcorner, x} \lr \underbrace{\vp(u) \text{ defines } x}_{(\star)}}
\end{align*}
Otherwise, the least undefinable ordinal would exist and be definable (because there are only countably many Gödel numbers).

Let's now go back to cardinals.

\begin{boxtheorem} \label{Ch3:Thm:H_Kappa_Sigma1}
    If $\kappa$ is an uncountable cardinal, then $H_{\kappa} \prec_{\Sigma_1} V$.
\end{boxtheorem}

Note that $H_{\kappa} \prec_{\Sigma_1} V$ means $\Sigma_1$ formulas are absolute between $H_{\kappa}$ and $V$. We borrow model theoretic notation for being an elementary substructure, but note that this is the model theory we know in the ``outside world'', not the model theory we have inside of $V$.

\begin{proof}
    Suppose $\vp(u)$ is the formula $\exists v\parenth{\psi\of{u, v}}$, with $\psi$ $\Delta_0$. Suppose $x \in H_{\kappa}$ is so that $\vp\of{x}$.

    We need to show that
    \begin{align*}
        \parenth{\exists v\parenth{\psi\of{x, v}}}^{H_{\kappa}}
    \end{align*}

    Pick any $y$ such that $\psi\of{x, y}$. Let $N = \trcl{\set{x, y}}$. Then $x, y \in N$ and $N$ is transitive. By $\Delta_0$ absoluteness, we know $\psi(x, y)^N$. Can write ``$N \models \psi(x, y)$''.

    % [FILLED] the proof from here on is supplied; the lecture broke off at the application of Downwards Löwenheim-Skolem
    Apply Downwards Löwenheim-Skolem to find $Z \prec N$ with $\trcl{\set{x}} \ssq Z$ and $\abs{Z} \leq \abs{\trcl{\set{x}}} + \aleph_0$. As $x \in H_{\kappa}$ and $\kappa$ is uncountable, $\abs{Z} < \kappa$.

    $N$ is transitive, so Extensionality holds in $N$, and hence in $Z \prec N$; so $\in$ is extensional on $Z$. It is also well-founded and set-like there, so by the Mostowski Collapse Theorem Scheme (\Cref{Ch2:Thm:Mostowski}) there is a transitive set $M$ and an isomorphism $\pi : \parenth{Z, \in} \to \parenth{M, \in}$, given by $\pi(z) = \setst{\pi(w)}{w \in z \cap Z}$. Two things about $\pi$:
    \begin{itemize}
        \item $\pi(w) = w$ for every $w \in \trcl{\set{x}}$. Indeed, by $\in$-induction: such a $w$ is a subset of the transitive set $\trcl{\set{x}} \ssq Z$, so $w \cap Z = w$, and $\pi(w) = \setst{\pi(v)}{v \in w} = w$. In particular, $\pi(x) = x$.
        \item $\abs{M} = \abs{Z} < \kappa$.
    \end{itemize}

    Now $N \models \exists v \, \psi\of{x, v}$, so, as $Z \prec N$ and $x \in Z$, there is $y' \in Z$ with $Z \models \psi\of{x, y'}$. Applying the isomorphism, $M \models \psi\of{x, \pi(y')}$, ie, $\psi\of{x, \pi(y')}^M$. As $M$ is transitive and $\psi$ is $\Delta_0$, $\Delta_0$-absoluteness gives $\psi\of{x, \pi(y')}$.

    Finally, $\pi(y') \in M$ and $M$ is transitive, so $\trcl{\set{\pi(y')}} \ssq M$ and $\abs{\trcl{\set{\pi(y')}}} \leq \abs{M} < \kappa$. So $\pi(y') \in H_{\kappa}$, and as $H_{\kappa}$ is transitive, $\Delta_0$-absoluteness gives $\psi\of{x, \pi(y')}^{H_{\kappa}}$. So $\parenth{\exists v \, \psi\of{x, v}}^{H_{\kappa}}$.
\end{proof}

% 5 Oct

We have see that $H_{\kappa} \prec_{\Sigma_1} V$ for all uncountable cardinals $\kappa$. It follows that for such $\kappa$, $\kappa = \beth_{\kappa}$ iff $H_{\kappa} = V_{\kappa}$, which in turn implies that $V_{\kappa} \prec_{\Sigma_1} V$.

Observe that $V_{\kappa} \prec_{\Sigma_1} V$ is \underline{not} a formula in the language of set theory. But $\kappa = \beth_{\kappa}$ and $H_{\kappa} = V_{\kappa}$ \textit{are}. By comprehension, put
\begin{align*}
    C_1 = \setst{\kappa}{\kappa = \beth_{\kappa}} = \setst{\kappa}{\kappa \text{ is uncountable and } H_{\kappa} = V_{\kappa}}
\end{align*}
Then $C_1$ is a class. Writing $C_1 = \setst{\kappa}{V_{\kappa} \prec_{\Sigma_1} V}$ is problematic.

Note that $C_1$ is club in $\OR$. Indeed, we proved that the set of $\beth$-fixed points contains a club class of ordinals. Define a class function $f : \OR \to \OR$ by
\begin{itemize}
    \item $f(0) = \beth_0$
    \item $f(\alpha + 1) = \beth_{f(\alpha) + 1}$
    \item $f(\beta) = \sup_{\alpha < \beta} f(\alpha)$
\end{itemize}
% [CORRECTED] was: ``for every limit $\beta \in \lim\of{\range\of{f}}$, then $\beta = f(\beta) = \beth_{\beta}$'' --- $f$ is strictly increasing, so $f(\beta) > \beta$ in general (eg, $\beta = f(\omega)$); what holds is that the values of $f$ at limits, which are exactly the limit points of its range, are $\beth$-fixed points
The range of $f$ is a cub and, for every limit ordinal $\beta$, $f(\beta) = \beth_{f(\beta)}$; that is, every element of $\lim\of{\range\of{f}}$\footnote{This is the class of limit points of $\range(f)$, and as we know, the class of limit points of a club is a club} is a $\beth$-fixed point. So $C_1$ contains a club. So $C_1$ is unbounded. Moreover, $C_1$ is closed. So it's a club in $\OR$.

\begin{boxtheorem}
    Suppose $\mu$ is a strongly inaccessible cardinal. Let $C = \setst{\kappa < \mu}{V_{\kappa} \prec V_{\mu}}$. Then $C$ os a club subset of $\mu$.
\end{boxtheorem}

\begin{remark}
    $V_{\kappa} \prec V_{\mu}$ \textbf{is} a formula in the language of set theory because both $V_{\kappa}$ and $V_{\mu}$ are \underline{sets}. It would be a different story if they were classes.

    ``This is model theory. Model theory is math. And math is embedded in set theory. So this makes sense!''
\end{remark}

\begin{proof}
    Clearly $C$ is closed in $\mu$. If $C \cap \lambda \ne \emptyset$ and $\lambda = \sup(C \cap \lambda)$, then we need $\lambda \in C$ (meaning of closed).

    For such a $\lambda$, we have $V_{\lambda} = \bigcup_{\kappa < \lambda} V_{\kappa} = \bigcup_{\kappa \in C \cap \lambda} V_{\kappa}$.

    \begin{itemize}
        \item It follows, for $\kappa < \kappa' < \mu$ with $\kappa, \kappa' \in C$, that $V_{\kappa} \prec V_{\kappa'}$.

        \item Then, $\grpres{V_{\kappa}}{\kappa \in C \cap \lambda}$ is a chain of elementary substructures of $V_{\mu}$, so its union is also an elementary substructure of $V_{\mu}$.

        \item $V_{\lambda} \prec V_{\mu}$, so $\lambda \in C$.

        \item It also follows that $V_{\kappa} \prec V_{\lambda}$ for $\kappa \in C \cap \lambda$.
    \end{itemize}

    Now let $\kappa_0 < \mu$ be an arbitrary uncountable cardinal. By Downwards Löwenheim-Skolem, we have $X_0 \prec V_{\mu}$ with $V_{\kappa_0} \cup \set{\kappa_0} \ssq X_0$ and $\abs{X_0} = \abs{V_{\kappa_0} \cup \set{\kappa_0}} = \beth_{\kappa_0} < \beth_{\mu} = \mu$.

    % IDEA: Take X_0, fill in things of smaller rank, then take X_1, fill in evrything of smaller rank, and so on. Each X_i is obtained by applying DLS to $V_{\mu}$. We then get $V_{\kappa} \bigcup_{n < \omega} V_{\kappa_n} = \bigcup_{n < \omega} X_n \prec V_{\kappa}$ where we konw the chain is elementary by the elementary chain theorem.
\end{proof}

% 7 Oct

\begin{boxdefinition}[Hierarchy]
    We say that a class function $\grpres{Z_{\alpha}}{\alpha \in C}$ is a \textbf{hierarchy} of sets iff
    \begin{itemize}
        \item each $Z_{\alpha}$ is a transitive set
        \item $C$ is a club class of ordinals
        \item for every $\alpha < \beta$ from $C$, $Z_{\alpha} \subset Z_{\beta}$
        \item for every $\beta$ a limit point\footnote{as $C$ is club, such $\beta$s do also lie in $C$} of $C$,
        \begin{align*}
            Z_{\beta} = \bigcup_{\alpha \in C \cap \beta} Z_{\alpha}
        \end{align*}
    \end{itemize}
    We usually write
    \begin{align*}
        Z = \bigcup_{\alpha \in C} Z_{\alpha}
        = \setst{x}{\exists \alpha \in C \st x \in Z_{\alpha}}
    \end{align*}
    (unless, of course, we already have a name for our hierarchy).
\end{boxdefinition}

\begin{boxexample}
    The following are all examples of hierarchies we've seen:
    \begin{align*}
        \grpres{V_{\alpha}}{\alpha \in \OR}
        \qquad , \qquad
        \grpres{H_{\kappa}}{\kappa \text{ an infinite cardinal}}
    \end{align*}
\end{boxexample}

Indeed, there will be examples where $Z \ne V$.

\begin{boxtheorem}[Reflection theorem scheme, proved in $\ZF$]
    Suppose $\varphi_i$ for $i < \ell < \omega$ are formulae in the language of set theory. Suppose that $\grpres{Z_{\alpha}}{\alpha \in C}$ is a hierarchy. Then there's a club class $D \ssq C$ such that for every $\alpha \in D$ and $i < \ell$, $\vp_i$ is absolute between $Z_{\alpha}$ and $Z$.
\end{boxtheorem}
\begin{proof}
    WLOG, assume that for all $j < \ell$ and proper subformulae $\psi$ of $\varphi_j$, there is some $i < j$ such that $\psi$ is $\vp_{i}$. Ie, we are setting up an induction on the complexity of $\vp$ by setting up an induction on ordinals.

    Let's define a class sequence of ordinals $\grpres{\alpha_{\eta}}{\eta \in \OR}$ as follows.
    \begin{itemize}
        % [CORRECTED] was: $\alpha_0 = 0$ --- $Z_{\alpha}$ is only defined for $\alpha \in C$, and $0$ need not be in $C$ (it is not for $\grpres{H_{\kappa}}{\kappa \text{ an infinite cardinal}}$), while the successor step uses $Z_{\alpha_0}$; the board also says to maintain $\alpha_{\eta} \in C$ for all $\eta$
        \item \underline{Base:} Define $\alpha_0 = \min(C)$.
        \item \underline{Successor:} Given $\alpha_{\eta}$, define $\alpha_{\eta + 1}$ to be the least $\beta \in C \setminus \parenth{\alpha_{\eta} + 1}$ so that for all $j < \ell$, if $\vp_{j}\of{\bar{u}}$ is of the form $\exists v, \vp_i\of{\bar{u}, v}$, then for all $\bar{x}$ from $Z_{\alpha_{\eta}}$ that match $\bar{u}$, if there is $y \in Z$ such that $\vp_i\of{\bar{x}, y}^Z$, then there is $y \in Z_{\beta}$ such that $\vp_i\of{\bar{x}, y}^Z$.
        \item \underline{Limit case:} If $\theta$ is a limit ordinal, define
        \begin{align*}
            \alpha_{\theta} = \sup_{\eta < \theta} \of{\alpha_{\eta}} \in \lim(C) \subset C
        \end{align*}
        (because $C$ is club).
    \end{itemize}

    The successor step is where the work happens. A tuple $\bar{x}$ from $Z_{\alpha_{\eta}}$ may have a witness $y'$ only higher up, in some $Z_{\beta'}$, and $\alpha_{\eta + 1}$ is chosen high enough to catch a witness for every such $\bar{x}$ at once:
    \begin{figure}[H]
        \centering
        \begin{tikzpicture}[dot/.style = {circle, fill, inner sep = 1.5pt}, every node/.style = {font = \small}]
            \drawcone{2.2}{5}{$Z$}
            \drawconelevel{2.2}{5}{1.6}{$Z_{\alpha_{\eta}}$}
            \drawconelevel{2.2}{5}{4.2}{$Z_{\beta'}$}
            \node[dot, label = {right:$\bar{x}$}] at (-0.15, 1.0) {};
            \node[dot, label = {right:$y'$}] at (0.3, 3.6) {};
        \end{tikzpicture}
    \end{figure}
    Only existential formulae need this treatment. We may take every $\vp_j$ to be built from atomic formulae using $\neg$, $\land$ and $\exists$ alone, writing $\forall$ as $\neg \exists \neg$. Relativization leaves atomic formulae unchanged, so they mean the same in $Z_{\alpha}$ as in $Z$, and it commutes with $\neg$ and $\land$, so absoluteness passes through them. The one place it can fail is $\exists v$, where $Z$ may have a witness that $Z_{\alpha}$ lacks.

    The successor step makes sense. For each $j < \ell$ of the form above and each $\bar{x}$ from $Z_{\alpha_{\eta}}$ that has a witness $y \in Z$, there is a least $\gamma \in C$ with a witness in $Z_{\gamma}$, as $Z = \bigcup_{\gamma \in C} Z_{\gamma}$. There are only set many such pairs $\parenth{j, \bar{x}}$, as $Z_{\alpha_{\eta}}$ is a set, so these $\gamma$ have a supremum, and as $C$ is unbounded there is some $\beta \in C$ above it and above $\alpha_{\eta}$. Any such $\beta$ works, since $Z_{\gamma} \ssq Z_{\beta}$ for $\gamma \leq \beta$ in $C$.

    Then take $D$ to be $\setst{\alpha_{\theta}}{\theta \text{ is a limit ordinal}}$. The sequence $\grpres{\alpha_{\eta}}{\eta \in \OR}$ is strictly increasing and continuous at limits, so $D$ is a club class: it is unbounded since $\alpha_{\theta} \geq \theta$, and closed since a supremum of $\alpha_{\theta}$s over limit ordinals $\theta$ in some set $\Theta$ with no largest element is $\alpha_{\sup \Theta}$, and $\sup \Theta$ is a limit ordinal. And $D \ssq C$ by the limit case.

    Finally, fix $\alpha = \alpha_{\theta} \in D$; we show that each $\vp_i$ is absolute between $Z_{\alpha}$ and $Z$, by induction on $i$. The atomic and propositional cases are as above, so suppose $\vp_j\of{\bar{u}}$ is $\exists v \, \vp_i\of{\bar{u}, v}$ with $i < j$, and let $\bar{x}$ from $Z_{\alpha}$ match $\bar{u}$.
    \begin{description}
        \item[\underline{From $Z_{\alpha}$ to $Z$.}] If $y \in Z_{\alpha}$ and $\vp_i\of{\bar{x}, y}^{Z_{\alpha}}$, then $y \in Z$ and $\vp_i\of{\bar{x}, y}^Z$ by the induction hypothesis, so $\vp_j\of{\bar{x}}^Z$.

        \item[\underline{From $Z$ to $Z_{\alpha}$.}] Suppose $\vp_j\of{\bar{x}}^Z$. As $\alpha_{\theta}$ is a limit point of $C$, $Z_{\alpha_{\theta}} = \bigcup_{\gamma \in C \cap \alpha_{\theta}} Z_{\gamma} = \bigcup_{\eta < \theta} Z_{\alpha_{\eta}}$, and $\bar{x}$ is a finite tuple, so $\bar{x}$ is from $Z_{\alpha_{\eta}}$ for some $\eta < \theta$. By the choice of $\alpha_{\eta + 1}$, there is $y \in Z_{\alpha_{\eta + 1}} \ssq Z_{\alpha}$ with $\vp_i\of{\bar{x}, y}^Z$, and by the induction hypothesis $\vp_i\of{\bar{x}, y}^{Z_{\alpha}}$. So $\vp_j\of{\bar{x}}^{Z_{\alpha}}$.
    \end{description}
\end{proof}

Here's an interesting corollary.

Let $\sigma_0, \ldots, \sigma_{\ell}$ be finitely many axioms of $\ZFC$. By the reflection theorem scheme, there are club many $\alpha \in \OR$ so that $\forall  i \leq \ell$, $\sigma_i^{V_{\alpha}}$. Equivalently, $\forall i \leq \ell$, $V_{\alpha} \models \sigma_i$. Rephrasing:

\begin{boxcorollary}
    For every finite sub=theory $T$ of $\ZFC$,
    \begin{align*}
        \ZFC \vdash \text{``$T$ has a transitive model''}
    \end{align*}
    where the quotes (``read with a different tone of voice'')
\end{boxcorollary}

With Choice, we can make the model countable as well. Given $\alpha$ with
\begin{align*}
    V_{\alpha} \models \bigwedge_{i \leq \ell} \sigma_i
\end{align*}
we can apply Downwards Löwenheim-Skolem to get some countable $X \prec V_{\alpha}$, so that $X \models \bigwedge_{i \leq \ell} \sigma_i$ too. As $V_{\alpha}$ is transitive, it satisfies Extensionality, and hence so does $X$; so $\in$ is extensional on $X$, as well as well-founded and set-like. Then apply Mostowski (\Cref{Ch2:Thm:Mostowski}) to get $X \simeq M$ with $M$ transitive and countable, and $M \models \bigwedge_{i \leq \ell} \sigma_i$, because satisfaction is preserved by isomorphisms. So in fact $\ZFC$ proves that every finite sub-theory of $\ZFC$ has a countable transitive model.





